\begin{proof}[Proof of \cref{obj:thm:finite-certificate}]
Fix the public parameters \((\beta,n,\sigma)\) satisfying the stated hypotheses. We first establish the three fixed-index guarantees simultaneously for every integer \(J\ge1\) and every \(P'\in\mathcal M_\beta(\sigma)\). For such \(J\) and \(P'\), define
\[
\begin{aligned}
A_{d,J}(P')
&=\int_0^1 K_J(x)f(P',s_dx)\mu_d(P',s_dx)\,dx,\\
B_{d,J}(P')
&=\int_0^1 K_J(x)f(P',s_dx)\,dx,\\
\eta_{d,J}(P')
&=\frac{A_{d,J}(P')}{B_{d,J}(P')},
\qquad d\in\{0,1\}.
\end{aligned}
\]
Here division by zero is assigned value zero; the positive denominator bound established below ensures that these are ordinary ratios throughout the model class. The signs \(s_1=1\) and \(s_0=-1\) are those of \cref{obj:synth_5}.

\begin{enumerate}
\item \textbf{Population identities and empirical moments.}
Define the single-observation summands, as functions on
\(\{0,1\}^2\times\mathbb R\), by
\[
\begin{aligned}
\xi^{\mathrm A}_{d,J}(D,Y,W)
&=\mathbf1\{D=d\}YQ_{J,\sigma}(s_dW),\\
\xi^{\mathrm B}_{d,J}(D,Y,W)
&=\mathbf1\{D=d\}Q_{J,\sigma}(s_dW).
\end{aligned}
\]
The inverse-heat polynomial is defined in \cref{obj:def:inverse-heat}.
Its evaluation at the proxy is square integrable: the latent score is bounded, Gaussian variables have every polynomial moment, and the indicator and binary outcome are bounded. Consequently,
\[
\xi^{\mathrm A}_{d,J},\xi^{\mathrm B}_{d,J}
\in L^2(P'_{\mathrm{obs}}).
\]

By \cref{obj:ass:gaussian,obj:ass:error-independence}, the reflected error \(s_d\varepsilon\) is standard Gaussian and independent of \((X,Y^0,Y^1)\). The first identity of \cref{obj:lem:exact-covariance} therefore gives, for every real \(x\),
\[
\int_{\mathbb R}Q_{J,\sigma}(s_dx+\sigma e)\,N(0,1)(de)
=K_J(s_dx).
\]
Integrating the independent Gaussian coordinate first, and using
\(\mathbf1\{D=d\}Y=\mathbf1\{D=d\}Y^d\), yields
\[
\begin{aligned}
\int \xi^{\mathrm B}_{d,J}\,dP'_{\mathrm{obs}}
&=\int \mathbf1\{D=d\}K_J(s_dX)\,dP',\\
\int \xi^{\mathrm A}_{d,J}\,dP'_{\mathrm{obs}}
&=\int \mathbf1\{D=d\}Y^dK_J(s_dX)\,dP'.
\end{aligned}
\]
The score marginal has density \(f(P',\cdot)\), so
\[
P'\{X=0\}=\int_{\{0\}}f(P',x)\,dx=0.
\]
Thus the arm indicator agrees almost surely with
\(\mathbf1\{s_dX\in[0,1]\}\). The defining conditional-mean identity gives the score marginal weighted by \(Y^d\) the density
\(\mu_d(P',x)f(P',x)\). Integrating the bounded arm weight against these two marginal identities gives
\[
\begin{aligned}
\int \xi^{\mathrm B}_{d,J}\,dP'_{\mathrm{obs}}
&=\int_{\{x:s_dx\in[0,1]\}}K_J(s_dx)f(P',x)\,dx
=B_{d,J}(P'),\\
\int \xi^{\mathrm A}_{d,J}\,dP'_{\mathrm{obs}}
&=\int_{\{x:s_dx\in[0,1]\}}
K_J(s_dx)\mu_d(P',x)f(P',x)\,dx
=A_{d,J}(P').
\end{aligned}
\]
For \(d=1\) the integration interval is \([0,1]\); for \(d=0\) it is \([-1,0]\), and substitution \(x\mapsto-x\) gives the displayed endpoint integrals.

For the second moments, the second identity of
\cref{obj:lem:exact-covariance}, integrated in the same order, gives
\[
\begin{aligned}
\int (\xi^{\mathrm B}_{d,J})^2\,dP'_{\mathrm{obs}}
&=\int_0^1 f(P',s_dx)
\sum_{j=0}^{m_J}\frac{\sigma^{2j}}{j!}
[K_J^{(j)}(x)]^2\,dx\\
&\le\frac34\int_0^1
\sum_{j=0}^{m_J}\frac{\sigma^{2j}}{j!}
[K_J^{(j)}(x)]^2\,dx
=V_J(\sigma),
\end{aligned}
\]
where the inequality uses \cref{obj:ass:density-bounds} and the nonnegativity of every summand. Since \(Y\in\{0,1\}\),
\[
(\xi^{\mathrm A}_{d,J})^2
=Y^2(\xi^{\mathrm B}_{d,J})^2
\le(\xi^{\mathrm B}_{d,J})^2,
\qquad
\int(\xi^{\mathrm A}_{d,J})^2\,dP'_{\mathrm{obs}}
\le V_J(\sigma).
\]

Define the centered empirical errors on the sample and seed space by
\[
\begin{aligned}
\Delta^{\mathrm A}_{d,J}(S_n,U)
&=\widehat A_{d,J}(S_n)-A_{d,J}(P'),\\
\Delta^{\mathrm B}_{d,J}(S_n,U)
&=\widehat B_{d,J}(S_n)-B_{d,J}(P').
\end{aligned}
\]
The empirical averages are those of \cref{obj:def:ratio}. For either
\(\diamond\in\{\mathrm A,\mathrm B\}\), the population identities show that
\[
\Delta^\diamond_{d,J}
=\frac1n\sum_{i=1}^n
\left[
\xi^\diamond_{d,J}(D_i,Y_i,W_i)
-\int\xi^\diamond_{d,J}\,dP'_{\mathrm{obs}}
\right].
\]
Independence makes the cross terms between distinct centered observations vanish. Expanding the square therefore gives
\[
\begin{aligned}
\mathbb E_{P',n}\Delta^\diamond_{d,J}&=0,\\
\mathbb E_{P',n}(\Delta^\diamond_{d,J})^2
&=\frac1n\int
\left(\xi^\diamond_{d,J}
-\int\xi^\diamond_{d,J}\,dP'_{\mathrm{obs}}\right)^2
dP'_{\mathrm{obs}}\\
&=\frac1n\left[
\int(\xi^\diamond_{d,J})^2\,dP'_{\mathrm{obs}}
-\left(\int\xi^\diamond_{d,J}\,dP'_{\mathrm{obs}}\right)^2
\right]\\
&\le\frac{V_J(\sigma)}n.
\end{aligned}
\]
These errors are square integrable. Integrating the independent seed contributes a factor one, since the averages depend on the sample. The same mean calculation proves
\[
\mathbb E_{P',n}\widehat B_{d,J}=B_{d,J}(P').
\]
% lean: CausalSmith.Stat.RdTruesideNoiseFrontier.marked_observed_integral_eq_endpoint
% lean: CausalSmith.Stat.RdTruesideNoiseFrontier.unmarked_observed_integral_eq_endpoint
% lean: CausalSmith.Stat.RdTruesideNoiseFrontier.marked_observed_sq_le_kernelVariance
% lean: CausalSmith.Stat.RdTruesideNoiseFrontier.unmarked_observed_sq_le_kernelVariance
% lean: CausalSmith.Stat.RdTruesideNoiseFrontier.seeded_iid_centered_moments
% lean: CausalSmith.Stat.RdTruesideNoiseFrontier.unmarkedMean_eq_single_observation

\item \textbf{Positive denominators and population bias.}
By \cref{obj:lem:positive-endpoint}, the kernel is nonnegative on \([0,1]\) and integrates to one. Hence \cref{obj:ass:density-bounds} gives
\[
\frac14
=\frac14\int_0^1K_J(x)\,dx
\le B_{d,J}(P')
\le\frac34\int_0^1K_J(x)\,dx
=\frac34.
\]
In particular, the denominator is positive and
\[
A_{d,J}(P')=\eta_{d,J}(P')B_{d,J}(P').
\]
The nonnegative weight \(K_J(x)f(P',s_dx)\), together with
\cref{obj:ass:mean-bounds}, gives
\[
\frac14B_{d,J}(P')
\le A_{d,J}(P')
\le\frac34B_{d,J}(P'),
\qquad
\eta_{d,J}(P')\in[1/4,3/4].
\]

For \(x\in[0,1]\), reflection preserves the distance from zero. Thus
\cref{obj:ass:mean-holder} gives
\[
|\mu_d(P',s_dx)-\mu_d(P',0)|\le2x^\beta.
\]
Subtracting the cutoff mean inside the weighted integral yields
\[
\eta_{d,J}(P')-\mu_d(P',0)
=
\frac{
\int_0^1K_J(x)f(P',s_dx)
[\mu_d(P',s_dx)-\mu_d(P',0)]\,dx
}{B_{d,J}(P')}.
\]
The numerator satisfies
\[
\begin{aligned}
\left|
\int_0^1K_J(x)f(P',s_dx)
[\mu_d(P',s_dx)-\mu_d(P',0)]\,dx
\right|
&\le\int_0^1K_J(x)f(P',s_dx)\,2x^\beta\,dx\\
&\le\frac32M_{J,\beta}.
\end{aligned}
\]
Since \(M_{J,\beta}\ge0\) and \(B_{d,J}(P')\ge1/4\), it follows that
\[
|\eta_{d,J}(P')-\mu_d(P',0)|\le6M_{J,\beta}.
\]
Together with the expectation identity above, this also proves the required lower bound on each unmarked population expectation.
% lean: CausalSmith.Stat.RdTruesideNoiseFrontier.endpointKernel_weighted_integral_bounds
% lean: CausalSmith.Stat.RdTruesideNoiseFrontier.arm_kernel_mean_bounds_and_bias

\item \textbf{Absolute risk.}
For every sample, define the floored denominator
\[
\widehat B^{\mathrm{floor}}_{d,J}
=\max\{\widehat B_{d,J},1/8\}.
\]
If \(\widehat B_{d,J}\ge1/8\), flooring leaves it unchanged. If
\(\widehat B_{d,J}<1/8\), then \(B_{d,J}(P')\ge1/4\) gives
\[
\left|\widehat B^{\mathrm{floor}}_{d,J}-B_{d,J}(P')\right|
=B_{d,J}(P')-\frac18
\le B_{d,J}(P')-\widehat B_{d,J}
=\left|\widehat B_{d,J}-B_{d,J}(P')\right|.
\]
Consequently, in both cases,
\[
\left|\widehat B^{\mathrm{floor}}_{d,J}-B_{d,J}(P')\right|
\le|\Delta^{\mathrm B}_{d,J}|,
\qquad
\widehat B^{\mathrm{floor}}_{d,J}\ge\frac18.
\]

Projection onto \([1/4,3/4]\) decreases distance to
\(\eta_{d,J}(P')\), which belongs to that interval. Using the population product identity,
\[
\frac{\widehat A_{d,J}}{\widehat B^{\mathrm{floor}}_{d,J}}
-\eta_{d,J}(P')
=
\frac{
\Delta^{\mathrm A}_{d,J}
+\eta_{d,J}(P')
[B_{d,J}(P')-\widehat B^{\mathrm{floor}}_{d,J}]
}{\widehat B^{\mathrm{floor}}_{d,J}}.
\]
Since \(|\eta_{d,J}(P')|\le3/4\), the projected estimate of
\cref{obj:def:ratio} satisfies
\[
|\widehat\mu_{d,J}-\eta_{d,J}(P')|
\le8|\Delta^{\mathrm A}_{d,J}|
+6|\Delta^{\mathrm B}_{d,J}|.
\]
This inequality applies to every sample, including those with an empty arm.

The triangle inequality for the two arm estimates and their population biases now gives
\[
\begin{aligned}
|\widehat\theta_J-\theta(P')|
\le{}&12M_{J,\beta}\\
&+8\bigl(|\Delta^{\mathrm A}_{1,J}|
+|\Delta^{\mathrm A}_{0,J}|\bigr)
+6\bigl(|\Delta^{\mathrm B}_{1,J}|
+|\Delta^{\mathrm B}_{0,J}|\bigr).
\end{aligned}
\]
For each of the four errors, Cauchy--Schwarz on the probability experiment gives
\[
\mathbb E_{P',n}|\Delta^\diamond_{d,J}|
\le\sqrt{\mathbb E_{P',n}(\Delta^\diamond_{d,J})^2}
\le\sqrt{\frac{V_J(\sigma)}n}.
\]
The displayed upper bound on the effect error is integrable, so integration yields
\[
\mathbb E_{P',n}|\widehat\theta_J-\theta(P')|
\le12M_{J,\beta}
+28\sqrt{\frac{V_J(\sigma)}n}.
\]
% lean: CausalSmith.Stat.RdTruesideNoiseFrontier.denominator_floor_error_le
% lean: CausalSmith.Stat.RdTruesideNoiseFrontier.clipped_ratio_error_le
% lean: CausalSmith.Stat.RdTruesideNoiseFrontier.thetaHat_absRisk_le_of_second_moments

\item \textbf{Fixed-index coverage.}
Kernel normalization implies
\[
\int_0^1K_J(x)^2\,dx>0.
\]
Indeed, if this nonnegative integral were zero, then \(K_J=0\) almost everywhere on \([0,1]\), contradicting
\(\int_0^1K_J(x)\,dx=1\). Every derivative term in \(V_J(\sigma)\) is nonnegative, and its order-zero coefficient is one, including at \(\sigma=0\). Hence
\[
V_J(\sigma)\ge\frac34\int_0^1K_J(x)^2\,dx>0.
\]
Define the positive threshold
\[
\zeta_J=\sqrt{\frac{40V_J(\sigma)}n}.
\]
The centered-moment identities imply, for every \(d\) and
\(\diamond\in\{\mathrm A,\mathrm B\}\),
\[
\operatorname{Var}_{P',n}(\Delta^\diamond_{d,J})
=\mathbb E_{P',n}(\Delta^\diamond_{d,J})^2
\le\frac{V_J(\sigma)}n.
\]
Chebyshev's inequality therefore gives
\[
\begin{aligned}
\mathbb P_{P',n}\{|\Delta^\diamond_{d,J}|>\zeta_J\}
&\le\mathbb P_{P',n}\{|\Delta^\diamond_{d,J}|\ge\zeta_J\}\\
&\le\frac{\operatorname{Var}_{P',n}(\Delta^\diamond_{d,J})}
{\zeta_J^2}
\le\frac1{40}.
\end{aligned}
\]
Define the simultaneous-control event on the sample and seed space by
\[
\Gamma_J
=\bigcap_{d\in\{0,1\}}
\left(
\{|\Delta^{\mathrm A}_{d,J}|\le\zeta_J\}
\cap
\{|\Delta^{\mathrm B}_{d,J}|\le\zeta_J\}
\right).
\]
The errors are measurable, and the union bound gives
\[
\mathbb P_{P',n}(\Gamma_J)
\ge1-\sum_{d\in\{0,1\}}
\sum_{\diamond\in\{\mathrm A,\mathrm B\}}
\mathbb P_{P',n}\{|\Delta^\diamond_{d,J}|>\zeta_J\}
\ge1-\frac4{40}=\frac9{10}.
\]
On this event the two arm errors relative to their weighted population means are each at most \(14\zeta_J\). Adding the two population biases gives
\[
|\widehat\theta_J-\theta(P')|
\le12M_{J,\beta}+28\zeta_J
=\rho_J,
\]
where \(\rho_J\) is the public radius of \cref{obj:synth_6}.

The cutoff mean bounds imply
\[
-\frac12\le\theta(P')\le\frac12.
\]
Thus, on \(\Gamma_J\),
\[
\max\{-1,\widehat\theta_J-\rho_J\}
\le\theta(P')
\le\min\{1,\widehat\theta_J+\rho_J\}.
\]
These are precisely the endpoints of \(I_J\) in
\cref{obj:synth_6}. Therefore
\[
\mathbb P_{P',n}\{\theta(P')\in I_J\}\ge\frac9{10}.
\]
Since \(J\ge1\) and \(P'\) were arbitrary, substituting \(J=L\) and \(P'=P\) proves all three fixed-index assertions.
% lean: CausalSmith.Stat.RdTruesideNoiseFrontier.kernelVariance_pos
% lean: CausalSmith.Stat.RdTruesideNoiseFrontier.centered_error_failure_le_one_fortieth
% lean: CausalSmith.Stat.RdTruesideNoiseFrontier.polyInterval_coverage_of_average_failure_bounds

\item \textbf{Honesty after public selection.}
By \cref{obj:def:public-selection}, \(L_\star\) is a deterministic function of the public parameters. Consider separately its zero and positive values.

If \(L_\star=0\), the candidate specified in \cref{obj:synth_6} is
\[
\widehat\theta_0=0,
\qquad
\rho_0=\frac12,
\qquad
I_0=[-1/2,1/2].
\]
The target bounds above give
\[
\mathbb P_{P',n}\{\theta(P')\in I_0\}=1
\qquad\text{for every }P'\in\mathcal M_\beta(\sigma).
\]
If \(L_\star\ge1\), the fixed-index coverage result gives
\[
\mathbb P_{P',n}\{\theta(P')\in I_{L_\star}\}
\ge\frac9{10}
\qquad\text{for every }P'\in\mathcal M_\beta(\sigma).
\]

For positive indices, projection of both arm estimates gives
\[
-\frac12\le\widehat\theta_J\le\frac12,
\qquad
\rho_J\ge0.
\]
Thus the centered interval and \([-1,1]\) both contain
\(\widehat\theta_J\); their intersection is a nonempty connected closed interval contained in \([-1,1]\). Its endpoints are measurable, being formed from polynomial averages, a positive denominator floor, projection, and finite maxima and minima. The zero-index interval has these properties as well. Since selection is deterministic, the selected interval has these same properties. Together with the uniform coverage bound, this proves
\[
I_{L_\star}\in\mathcal I_{\beta,n,\sigma}
\]
by \cref{obj:def:honest-class}.
% lean: CausalSmith.Stat.RdTruesideNoiseFrontier.polyInterval_properties
% lean: CausalSmith.Stat.RdTruesideNoiseFrontier.honestAttainer_of_positive_coverage

\item \textbf{Selected risk.}
If \(L_\star=0\), then for every admissible \(P'\),
\[
|\widehat\theta_0-\theta(P')|
=|\theta(P')|\le\frac12=\rho_0,
\qquad
\mathbb E_{P',n}|\widehat\theta_0-\theta(P')|\le\rho_0.
\]
If \(L_\star\ge1\), nonnegativity of \(V_{L_\star}(\sigma)/n\) and monotonicity of the square root give
\[
\begin{aligned}
\mathbb E_{P',n}|\widehat\theta_{L_\star}-\theta(P')|
&\le12M_{L_\star,\beta}
+28\sqrt{\frac{V_{L_\star}(\sigma)}n}\\
&\le12M_{L_\star,\beta}
+28\sqrt{\frac{40V_{L_\star}(\sigma)}n}
=\rho_{L_\star}.
\end{aligned}
\]
In both cases the bound is uniform over \(P'\). Since
\(E_\beta(n,\sigma)=\rho_{L_\star}\) by
\cref{obj:def:public-selection}, taking the supremum gives
\[
\sup_{P'\in\mathcal M_\beta(\sigma)}
\mathbb E_{P',n}|\widehat\theta_{L_\star}-\theta(P')|
\le E_\beta(n,\sigma).
\]
% lean: CausalSmith.Stat.RdTruesideNoiseFrontier.attainerRisk_of_positive_risk

\item \textbf{Selected expected length.}
For every positive candidate index \(J\), the endpoint formula gives samplewise
\[
\begin{aligned}
\operatorname{length}(I_J)
&=\max\!\left\{0,\,
\min\{1,\widehat\theta_J+\rho_J\}
-\max\{-1,\widehat\theta_J-\rho_J\}\right\}\\
&\le2\rho_J.
\end{aligned}
\]
Indeed, the upper endpoint is at most
\(\widehat\theta_J+\rho_J\), the lower endpoint is at least
\(\widehat\theta_J-\rho_J\), and \(2\rho_J\ge0\).
For the zero candidate,
\[
\operatorname{length}(I_0)=1=2\rho_0.
\]
Consequently the bound holds at the selected index for every sample. Integrating under the probability experiment and then taking the supremum yields
\[
\sup_{P'\in\mathcal M_\beta(\sigma)}
\mathbb E_{P',n}\operatorname{length}(I_{L_\star})
\le2\rho_{L_\star}
=2E_\beta(n,\sigma).
\]
% lean: CausalSmith.Stat.RdTruesideNoiseFrontier.attainerLength_le_certificate
\end{enumerate}
\end{proof}

\begin{proof}[Proof of \cref{obj:thm:uniform-frontier}]
Fix \(0<\beta\le1\). We use the Legendre identities of
\cref{obj:lem:classical-legendre-facts} and the Hermite generating function and orthogonality of
\cref{obj:lem:classical-hermite-facts}. All constants introduced below depend at most on \(\beta\), unless explicitly indexed by an additional fixed parameter.

\begin{enumerate}
\item \textbf{Root facts used in the comparisons.}
For \(n\ge2\) and \(0\le\sigma\le1\), use the notation of
\cref{obj:def:uniform-rate}:
\[
\nu=2\beta+1,\qquad h_0=n^{-1/\nu}.
\]
Define, on \(0<h\le1\),
\[
\gamma_{\mathrm{bal}}(h)
=\log n+\nu\log h-\Psi(h,\sigma).
\]
We have
\[
0<h_0\le1,\qquad
\log n+\nu\log h_0=0.
\]
By \cref{obj:lem:noise-cost-scaling}, the function
\(\gamma_{\mathrm{bal}}\) is continuous. For \(0<h<t\le1\),
\[
\gamma_{\mathrm{bal}}(t)-\gamma_{\mathrm{bal}}(h)
=\nu\log(t/h)+\Psi(h,\sigma)-\Psi(t,\sigma)>0.
\]
Its endpoint values satisfy
\[
\gamma_{\mathrm{bal}}(h_0)=-\Psi(h_0,\sigma)\le0,
\qquad
\gamma_{\mathrm{bal}}(1)=\log n>0.
\]
The intermediate value theorem therefore gives a unique zero
\(\xi_{\mathrm{root}}\in[h_0,1]\). Strict increase also gives
\[
\{h\in[h_0,1]:\gamma_{\mathrm{bal}}(h)\le0\}
=[h_0,\xi_{\mathrm{root}}].
\]
Consequently, the supremum of this sublevel set equals
\(\xi_{\mathrm{root}}\), and the resolution in
\cref{obj:def:uniform-rate} satisfies
\[
h_\star(n,\sigma)=\xi_{\mathrm{root}},
\qquad
\log n+\nu\log h_\star(n,\sigma)
=\Psi(h_\star(n,\sigma),\sigma).
\]
Exponentiating yields the balance identity
\[
\exp\{\Psi(h_\star(n,\sigma),\sigma)\}
=n\,h_\star(n,\sigma)^\nu.
\]

If \(\sigma\le h_0\), then \(\Psi(h_0,\sigma)=0\), so uniqueness gives
\[
h_\star(n,\sigma)=h_0,\qquad
r_\beta(n,\sigma)=n^{-\beta/\nu}.
\]
This includes \(\sigma=0\).

If \(\sigma>h_0\), then \(\Psi(h_0,\sigma)>0\). Indeed, when
\(\sigma^4\le h_0\),
\[
(\sigma/h_0)^{2/3}>1.
\]
When \(h_0<\sigma^4\), we have \(h_0<1\) and
\(\sqrt{h_0}<\sigma^2\), whence
\[
h_0^{-1/2}>1,\qquad
1+\log(\sigma^2/\sqrt{h_0})\ge1,
\]
and their product exceeds one. Thus
\[
\gamma_{\mathrm{bal}}(h_0)<0,
\qquad
\gamma_{\mathrm{bal}}(\sigma)
=\log n+\nu\log\sigma
>\log n+\nu\log h_0=0.
\]
Strict increase now gives
\[
h_0<h_\star(n,\sigma)<\sigma.
\]
% lean: CausalSmith.Stat.RdTruesideNoiseFrontier.unique_rate_root
% lean: CausalSmith.Stat.RdTruesideNoiseFrontier.rateRoot_exp_balance
% lean: CausalSmith.Stat.RdTruesideNoiseFrontier.rateResolution_eq_direct_of_le
% lean: CausalSmith.Stat.RdTruesideNoiseFrontier.noisy_rateResolution_interior

\item \textbf{Lower comparisons.}
By \cref{obj:thm:direct-reduction}, there is \(c_{\mathrm{dir}}>0\) such that
\[
R_\beta(n,\sigma)\ge c_{\mathrm{dir}}n^{-\beta/\nu},
\qquad
\mathcal L_\beta(n,\sigma)\ge c_{\mathrm{dir}}n^{-\beta/\nu}
\]
for every admissible \(n,\sigma\). These are the required lower comparisons when
\(\sigma\le h_0\).

Consider \(\sigma>h_0\). For each integer \(m\ge2\), the scales of
\cref{obj:def:noise-support} satisfy
\[
b_m(\sigma)=\min\{1,\sigma\sqrt m/8\},
\qquad
h_m(\sigma)=b_m(\sigma)/m^2,
\qquad
0<b_m(\sigma)\le1.
\]
The support scale is nondecreasing in \(m\). Since
\(m^2\le4(m-1)^2\) for \(m\ge2\), it follows that
\[
h_{m-1}(\sigma)
=\frac{b_{m-1}(\sigma)}{(m-1)^2}
\le\frac{b_m(\sigma)}{(m-1)^2}
\le4h_m(\sigma).
\]
Also \(h_m(\sigma)\le m^{-2}\). In particular, any integer
\[
m>\sqrt{\frac{4096}{h_\star(n,\sigma)}}+2
\]
satisfies \(m\ge2\) and \(h_m(\sigma)\le h_\star(n,\sigma)/4096\).
Choose the smallest integer \(m\ge2\) with this latter property. At index two,
\[
b_2(\sigma)=\frac{\sigma\sqrt2}{8},
\qquad
h_2(\sigma)=\frac{\sigma\sqrt2}{32}
>\frac{h_\star(n,\sigma)}{4096},
\]
because \(h_\star(n,\sigma)<\sigma\) and
\(1\le\sqrt2<2\). Hence \(m\ge3\), and minimality gives
\[
h_{m-1}(\sigma)>\frac{h_\star(n,\sigma)}{4096}.
\]
Combining this with the predecessor bound gives
\[
\frac{h_\star(n,\sigma)}{16384}
<h_m(\sigma)
\le\frac{h_\star(n,\sigma)}{4096}.
\]

For this chosen index, define
\[
h_{\mathrm{low}}=h_m(\sigma),\qquad
b_{\mathrm{low}}=b_m(\sigma),\qquad
F_{\mathrm{cost}}(h)=1+\Psi(h,\sigma)\quad(0<h\le1).
\]
Thus \(0<h_{\mathrm{low}}\le h_\star(n,\sigma)<\sigma\le1\).
The ratio inequality of \cref{obj:lem:noise-cost-scaling} gives
\[
\frac{F_{\mathrm{cost}}(h_{\mathrm{low}})}
     {F_{\mathrm{cost}}(h_\star(n,\sigma))}
\ge
\left(\frac{h_\star(n,\sigma)}{h_{\mathrm{low}}}\right)^{1/2}
\ge64.
\]
Using \(h_{\mathrm{low}}^\nu\le h_\star(n,\sigma)^\nu\), the balance identity, and
\(\Psi(h_\star(n,\sigma),\sigma)\ge0\), we obtain
\[
\begin{aligned}
n h_{\mathrm{low}}^\nu
 \exp\{-F_{\mathrm{cost}}(h_{\mathrm{low}})/32\}
&\le
n h_\star(n,\sigma)^\nu
 \exp\{-2F_{\mathrm{cost}}(h_\star(n,\sigma))\}\\
&=\exp\{-2-\Psi(h_\star(n,\sigma),\sigma)\}\\
&\le e^{-2}.
\end{aligned}
\]

Define the observed alternative laws and their information budget by
\[
\mathsf P_\pm=(P^\pm_{b_{\mathrm{low}},m})_{\mathrm{obs}},
\qquad
\Delta_{\mathrm{info}}
=\frac{32}{3}\kappa^2h_{\mathrm{low}}^\nu
 \exp\{-F_{\mathrm{cost}}(h_{\mathrm{low}})/32\},
\qquad
\kappa=\frac1{100}.
\]
By \cref{obj:lem:cancellation-information-envelope},
\[
\chi^2(\mathsf P_+,\mathsf P_-)\le\Delta_{\mathrm{info}}<\infty,
\]
and the preceding calculation gives
\[
0\le n\Delta_{\mathrm{info}}
\le\frac{32}{3}\kappa^2e^{-2}
\le\frac{32}{30000}
\le\frac1{100}.
\]
Finiteness of this divergence supplies absolute continuity and an integrable
squared likelihood-ratio deviation. To spell out the product comparison, put
\[
\chi_{\mathrm{obs}}=\chi^2(\mathsf P_+,\mathsf P_-),
\qquad
\chi_{\mathrm{sample}}
=\chi^2(\mathsf P_+^{\otimes n},\mathsf P_-^{\otimes n}).
\]
The product likelihood identity and \(1+\chi_{\mathrm{obs}}\le e^{\chi_{\mathrm{obs}}}\) give
\[
1+\chi_{\mathrm{sample}}
=(1+\chi_{\mathrm{obs}})^n
\le e^{n\chi_{\mathrm{obs}}}
\le e^{1/100}
\le\frac1{1-1/100}
=\frac{100}{99}.
\]
Therefore
\[
\chi_{\mathrm{sample}}\le\frac1{99}\le\frac4{25},
\qquad
\operatorname{TV}(\mathsf P_+^{\otimes n},\mathsf P_-^{\otimes n})
\le\frac12\sqrt{\chi_{\mathrm{sample}}}
\le\frac15.
\]
Here the total-variation inequality follows by applying Cauchy--Schwarz to
the likelihood-ratio deviation; the positive and negative parts of that
deviation have equal mass.

The support and order conditions required by
\cref{obj:lem:two-point-transfer} hold. That result, which includes independent
procedure randomization, yields
\[
R_\beta(n,\sigma)\ge\frac45\kappa h_{\mathrm{low}}^\beta,
\qquad
\mathcal L_\beta(n,\sigma)\ge\frac65\kappa h_{\mathrm{low}}^\beta.
\]
Consequently, with
\[
c_{\mathrm{noise}}=\frac45\kappa\,16384^{-\beta},
\qquad
c_{\mathrm{mat}}=\min\{c_{\mathrm{dir}},c_{\mathrm{noise}}\}>0,
\]
the direct and noisy cases together give
\[
c_{\mathrm{mat}}r_\beta(n,\sigma)\le R_\beta(n,\sigma),
\qquad
c_{\mathrm{mat}}r_\beta(n,\sigma)\le\mathcal L_\beta(n,\sigma)
\]
uniformly over \(n\ge2\) and \(0\le\sigma\le1\).
% lean: CausalSmith.Stat.RdTruesideNoiseFrontier.exists_noiseResolution_bracket
% lean: CausalSmith.Stat.RdTruesideNoiseFrontier.rounded_noise_information_factor
% lean: CausalSmith.Stat.RdTruesideNoiseFrontier.sampleLaw_tv_le_one_fifth_of_chiSqExt
% lean: CausalSmith.Stat.RdTruesideNoiseFrontier.frontierRate_lower_bounds

\item \textbf{Upper comparison for the public radius.}
Choose \(K_{\mathrm{bias}}>0\) and \(D_{\mathrm{env}}>0\) from
\cref{obj:lem:positive-endpoint,obj:lem:variance-cost-envelope}, so that
\[
M_{L,\beta}\le K_{\mathrm{bias}}L^{-2\beta},
\qquad
V_L(\sigma)\le D_{\mathrm{env}}L^2
 \exp\{D_{\mathrm{env}}[1+\Psi(L^{-2},\sigma)]\}
\]
for every \(L\ge1\). Define
\[
\begin{aligned}
A_{\mathrm{dil}}&=4(D_{\mathrm{env}}+1)^2,\\
Q_{\mathrm{var}}&=\sqrt{160D_{\mathrm{env}}e^{D_{\mathrm{env}}+1}},\\
C_{\mathrm{up}}&=\max\left\{
 \frac12A_{\mathrm{dil}}^\beta,\,
 12K_{\mathrm{bias}}A_{\mathrm{dil}}^\beta+28Q_{\mathrm{var}}
 \right\}.
\end{aligned}
\]
These constants are positive, and \(A_{\mathrm{dil}}\ge4\).

The radii in \cref{obj:synth_6} and the public minimization in
\cref{obj:def:public-selection} give
\[
\rho_0=\frac12,\qquad
\rho_L=12M_{L,\beta}+28\sqrt{\frac{40V_L(\sigma)}n}\quad(L\ge1),
\qquad
E_\beta(n,\sigma)\le\rho_L\quad(0\le L\le L_{\max}),
\]
where the finite cap is
\[
L_{\max}=\left\lceil n^{1/(4\beta+2)}\right\rceil.
\]

If \(1\le A_{\mathrm{dil}}h_\star(n,\sigma)\), then
\[
E_\beta(n,\sigma)\le\frac12
\le\frac12A_{\mathrm{dil}}^\beta h_\star(n,\sigma)^\beta
\le C_{\mathrm{up}}r_\beta(n,\sigma).
\]
Suppose instead that \(A_{\mathrm{dil}}h_\star(n,\sigma)<1\), and define
\[
x_{\mathrm{up}}
=(A_{\mathrm{dil}}h_\star(n,\sigma))^{-1/2},
\qquad
L_{\mathrm{up}}=\lceil x_{\mathrm{up}}\rceil,
\qquad
h_{\mathrm{up}}=L_{\mathrm{up}}^{-2}.
\]
Since \(x_{\mathrm{up}}>1\), ceiling bounds give
\[
1\le L_{\mathrm{up}},
\qquad
x_{\mathrm{up}}\le L_{\mathrm{up}}
<x_{\mathrm{up}}+1\le2x_{\mathrm{up}}.
\]
Taking inverse squares yields
\[
\frac{A_{\mathrm{dil}}h_\star(n,\sigma)}4
\le h_{\mathrm{up}}
\le A_{\mathrm{dil}}h_\star(n,\sigma)<1.
\]
Moreover,
\[
x_{\mathrm{up}}
\le h_\star(n,\sigma)^{-1/2}
\le h_0^{-1/2}
=n^{1/(4\beta+2)},
\]
so monotonicity of the ceiling gives \(L_{\mathrm{up}}\le L_{\max}\).

Define the dilation multiplier by
\[
\alpha_{\mathrm{dil}}=\frac{h_{\mathrm{up}}}{h_\star(n,\sigma)}.
\]
The width bounds imply
\[
(D_{\mathrm{env}}+1)^2\le\alpha_{\mathrm{dil}}\le A_{\mathrm{dil}},
\qquad
\alpha_{\mathrm{dil}}\ge1,
\qquad
\alpha_{\mathrm{dil}}h_\star(n,\sigma)=h_{\mathrm{up}}\le1.
\]
Thus \cref{obj:lem:noise-cost-scaling} applies and gives
\[
1+\Psi(h_{\mathrm{up}},\sigma)
\le\max\left\{
1,\,
\alpha_{\mathrm{dil}}^{-1/2}
[1+\Psi(h_\star(n,\sigma),\sigma)]
\right\}.
\]
Because \(D_{\mathrm{env}}/\sqrt{\alpha_{\mathrm{dil}}}\le1\), multiplication by
\(D_{\mathrm{env}}\) gives
\[
\begin{aligned}
D_{\mathrm{env}}[1+\Psi(h_{\mathrm{up}},\sigma)]
&\le
\max\left\{
D_{\mathrm{env}},\,
1+\Psi(h_\star(n,\sigma),\sigma)
\right\}\\
&\le D_{\mathrm{env}}+1+\Psi(h_\star(n,\sigma),\sigma).
\end{aligned}
\]
Hence
\[
V_{L_{\mathrm{up}}}(\sigma)
\le D_{\mathrm{env}}L_{\mathrm{up}}^2
 e^{D_{\mathrm{env}}+1}
 \exp\{\Psi(h_\star(n,\sigma),\sigma)\}.
\]
The rounding bound also gives
\[
L_{\mathrm{up}}^2
\le\frac4{A_{\mathrm{dil}}h_\star(n,\sigma)}
\le\frac4{h_\star(n,\sigma)}.
\]
Using the root balance and \(\nu=2\beta+1\), we obtain
\[
\frac{40V_{L_{\mathrm{up}}}(\sigma)}n
\le160D_{\mathrm{env}}e^{D_{\mathrm{env}}+1}
 h_\star(n,\sigma)^{2\beta},
\]
and therefore
\[
\sqrt{\frac{40V_{L_{\mathrm{up}}}(\sigma)}n}
\le Q_{\mathrm{var}}h_\star(n,\sigma)^\beta.
\]
The bias bound gives
\[
M_{L_{\mathrm{up}},\beta}
\le K_{\mathrm{bias}}h_{\mathrm{up}}^\beta
\le K_{\mathrm{bias}}A_{\mathrm{dil}}^\beta
 h_\star(n,\sigma)^\beta.
\]
Since \(L_{\mathrm{up}}\) is an eligible candidate,
\[
\begin{aligned}
E_\beta(n,\sigma)
&\le\rho_{L_{\mathrm{up}}}\\
&\le
(12K_{\mathrm{bias}}A_{\mathrm{dil}}^\beta+28Q_{\mathrm{var}})
 h_\star(n,\sigma)^\beta\\
&\le C_{\mathrm{up}}r_\beta(n,\sigma).
\end{aligned}
\]
Both cases establish this comparison over the full parameter range.
% lean: CausalSmith.Stat.RdTruesideNoiseFrontier.upperDegree_width_bounds
% lean: CausalSmith.Stat.RdTruesideNoiseFrontier.upperDegree_le_Lmax
% lean: CausalSmith.Stat.RdTruesideNoiseFrontier.dilation_cost_absorption
% lean: CausalSmith.Stat.RdTruesideNoiseFrontier.certificate_le_frontierRate

\item \textbf{Matched estimation and interval chains.}
By \cref{obj:thm:finite-certificate}, the selected procedures satisfy
\[
\begin{aligned}
\sup_{P\in\mathcal M_\beta(\sigma)}
\mathbb E_{P,n}
|\widetilde\theta_{\beta,n,\sigma}-\theta(P)|
&\le E_\beta(n,\sigma),\\
\sup_{P\in\mathcal M_\beta(\sigma)}
\mathbb E_{P,n}\operatorname{length}(\widetilde I_{\beta,n,\sigma})
&\le2E_\beta(n,\sigma),\\
\widetilde I_{\beta,n,\sigma}&\in\mathcal I_{\beta,n,\sigma}.
\end{aligned}
\]
The minimax infima are bounded above by the objectives of these eligible procedures.

For completeness, the mean bounds imply
\[
|\theta(P)|=|\mu_1(P,0)-\mu_0(P,0)|\le\frac12.
\]
Thus the constant estimator zero has worst expected absolute error at most
\(1/2\), and the constant interval \([-1,1]\) has coverage one and length two.
Both extended minimax objectives are consequently finite. Conversion to real
numbers preserves their comparisons with the selected procedures.

Put
\[
C_{\mathrm{mat}}=2C_{\mathrm{up}}>0.
\]
Combining the lower comparisons, the finite certificate, and the upper comparison gives
\[
\begin{aligned}
c_{\mathrm{mat}}r_\beta(n,\sigma)
&\le R_\beta(n,\sigma)
\le\sup_{P\in\mathcal M_\beta(\sigma)}
 \mathbb E_{P,n}|\widetilde\theta_{\beta,n,\sigma}-\theta(P)|\\
&\le E_\beta(n,\sigma)
\le C_{\mathrm{mat}}r_\beta(n,\sigma),\\
c_{\mathrm{mat}}r_\beta(n,\sigma)
&\le\mathcal L_\beta(n,\sigma)
\le\sup_{P\in\mathcal M_\beta(\sigma)}
 \mathbb E_{P,n}\operatorname{length}(\widetilde I_{\beta,n,\sigma})\\
&\le2E_\beta(n,\sigma)
\le C_{\mathrm{mat}}r_\beta(n,\sigma).
\end{aligned}
\]
The lower transfer ranges over the stated randomized decision classes, with
arbitrary procedure degree.
% lean: CausalSmith.Stat.RdTruesideNoiseFrontier.matchedFrontier_of_scalar_bounds
% lean: CausalSmith.Stat.RdTruesideNoiseFrontier.matched_frontier

\item \textbf{Finite-sample branches and interfaces.}
Define the branch comparison constants by
\[
\begin{gathered}
A_{\mathrm{int}}=(1+3\nu/2)^{3/2},\qquad
c_\tau=\frac1{2(1+2\nu)},\qquad C_\tau=e+6,\\
c_{\mathrm{pow}}=C_\tau^{-2\beta},\qquad
C_{\mathrm{pow}}=c_\tau^{-2\beta},\\
c_{\mathrm{br}}=\min\{1,c_\tau,c_{\mathrm{pow}}\},
\qquad
C_{\mathrm{br}}=\max\{A_{\mathrm{int}},C_\tau,C_{\mathrm{pow}}\}.
\end{gathered}
\]
All are positive.

The direct branch was established above. In the noisy region,
\(0<\sigma\le1\) and \(h_0<h_\star(n,\sigma)<\sigma\). Evaluation at the interface gives
\[
\Psi(\sigma^4,\sigma)=\sigma^{-2}-1,
\qquad
\gamma_{\mathrm{bal}}(\sigma^4)
=\log(n\sigma^{4\nu})-(\sigma^{-2}-1).
\]
This identity also holds at \(\sigma=1\), where both quantities on the right
of the noise-cost identity are zero. Strict increase on \((0,1]\) gives
\[
\sigma^4\le h_\star(n,\sigma)
\quad\Longleftrightarrow\quad
\log(n\sigma^{4\nu})\le\sigma^{-2}-1.
\]
The comparison is valid throughout the noisy region, including when
\(\sigma^4<h_0\).

On this intermediate branch, define
\[
z=\left(\frac{\sigma}{h_\star(n,\sigma)}\right)^{2/3},
\qquad
S_{\mathrm{int}}=1+\log(n\sigma^\nu).
\]
The inequalities \(h_\star(n,\sigma)<\sigma\) and
\(\sigma^4\le h_\star(n,\sigma)\) imply
\[
1<z\le\sigma^{-2}.
\]
The intermediate noise-cost formula and the root identity give
\[
1+\Psi(h_\star(n,\sigma),\sigma)=z,
\qquad
\log z=\frac23\{\log\sigma-\log h_\star(n,\sigma)\}.
\]
Substitution yields
\[
z+\frac{3\nu}{2}\log z=S_{\mathrm{int}},
\qquad
h_\star(n,\sigma)=\sigma z^{-3/2}.
\]
The left side of the scalar equation is strictly increasing for \(z>0\),
so the indicated root is unique. Since \(z>1\),
\[
0\le\log z\le z,
\qquad
z\le S_{\mathrm{int}}\le(1+3\nu/2)z.
\]
Taking inverse powers gives
\[
\frac{\sigma}{S_{\mathrm{int}}^{3/2}}
\le h_\star(n,\sigma)
\le A_{\mathrm{int}}\frac{\sigma}{S_{\mathrm{int}}^{3/2}}.
\]

If the branch-test inequality is strict in the other direction, then
\(h_\star(n,\sigma)<\sigma^4\). Define
\[
\tau=h_\star(n,\sigma)^{-1/2}.
\]
We have
\[
\tau>\sigma^{-2}\ge1,\qquad
h_\star(n,\sigma)=\tau^{-2},
\]
and the compact cost formula gives
\[
1+\Psi(h_\star(n,\sigma),\sigma)
=\tau\{1+\log(\sigma^2\tau)\}.
\]
Since \(\log\tau=-\tfrac12\log h_\star(n,\sigma)\), the root equation becomes
\[
2\nu\log\tau+\tau\{1+\log(\sigma^2\tau)\}
=1+\log n.
\]
For \(\sigma^{-2}<\tau_1<\tau_2\), both
\(\log\tau\) and \(1+\log(\sigma^2\tau)\) increase strictly, and the latter
factor is positive. In particular,
\[
\begin{aligned}
\tau_1\{1+\log(\sigma^2\tau_1)\}
&<\tau_2\{1+\log(\sigma^2\tau_1)\}\\
&<\tau_2\{1+\log(\sigma^2\tau_2)\}.
\end{aligned}
\]
Thus the compact equation has a unique root on the indicated domain.

To compare that root with an explicit expression, define
\[
\Lambda_n=1+\log n,\qquad
\Omega_{n,\sigma}=\log(e+\sigma^2\Lambda_n),\qquad
y_{\mathrm{cmp}}=\sigma^2\tau,\qquad
q_{\mathrm{cmp}}=\frac{\Lambda_n}{\tau},\qquad
J_{\mathrm{cmp}}=1+\log y_{\mathrm{cmp}}.
\]
Here \(\Lambda_n>1\), \(\Omega_{n,\sigma}>0\),
\(y_{\mathrm{cmp}}>1\), and \(J_{\mathrm{cmp}}>1\). Dividing the compact
equation by \(\tau\) gives
\[
q_{\mathrm{cmp}}
=2\nu\frac{\log\tau}{\tau}+J_{\mathrm{cmp}}.
\]
Using \(0\le\log\tau\le\tau\) and \(\nu\le3\), we obtain
\[
J_{\mathrm{cmp}}\le q_{\mathrm{cmp}}
\le J_{\mathrm{cmp}}+2\nu
\le J_{\mathrm{cmp}}+6,
\qquad
q_{\mathrm{cmp}}\le(1+2\nu)J_{\mathrm{cmp}}.
\]
Also \(\sigma^2\Lambda_n=y_{\mathrm{cmp}}q_{\mathrm{cmp}}\), so
\[
\Omega_{n,\sigma}\ge1,\qquad
\Omega_{n,\sigma}\ge\log y_{\mathrm{cmp}},
\qquad
J_{\mathrm{cmp}}\le2\Omega_{n,\sigma}.
\]
For the reverse comparison,
\[
e+\sigma^2\Lambda_n
=e+y_{\mathrm{cmp}}q_{\mathrm{cmp}}
\le y_{\mathrm{cmp}}(e+J_{\mathrm{cmp}}+6).
\]
The elementary bound \(\log t\le t-1\) for \(t>0\) therefore gives
\[
\begin{aligned}
\Omega_{n,\sigma}
&\le\log y_{\mathrm{cmp}}+\log(e+J_{\mathrm{cmp}}+6)\\
&\le(J_{\mathrm{cmp}}-1)+(e+J_{\mathrm{cmp}}+5)\\
&=2J_{\mathrm{cmp}}+e+4\\
&\le(e+6)J_{\mathrm{cmp}}.
\end{aligned}
\]
It follows that
\[
q_{\mathrm{cmp}}\le2(1+2\nu)\Omega_{n,\sigma},
\qquad
\Omega_{n,\sigma}\le(e+6)q_{\mathrm{cmp}}.
\]
All quantities divided below are positive, and these inequalities yield
\[
c_\tau\frac{\Lambda_n}{\Omega_{n,\sigma}}
\le\tau
\le C_\tau\frac{\Lambda_n}{\Omega_{n,\sigma}}.
\]
Raising this comparison to the negative power \(-2\beta\) reverses the
inequalities. Since
\[
r_\beta(n,\sigma)=\tau^{-2\beta},
\qquad
\left(\frac{\Lambda_n}{\Omega_{n,\sigma}}\right)^{-2\beta}
=\left(\frac{\Omega_{n,\sigma}}{\Lambda_n}\right)^{2\beta},
\]
we obtain
\[
c_{\mathrm{pow}}
\left(\frac{\Omega_{n,\sigma}}{\Lambda_n}\right)^{2\beta}
\le r_\beta(n,\sigma)
\le C_{\mathrm{pow}}
\left(\frac{\Omega_{n,\sigma}}{\Lambda_n}\right)^{2\beta}.
\]
These are the asserted compact comparisons because
\(\Lambda_n=\log(en)\).

At the intermediate--compact equality,
\[
\gamma_{\mathrm{bal}}(\sigma^4)=0.
\]
If \(\sigma^4<h_0\), strict increase would give
\(0=\gamma_{\mathrm{bal}}(\sigma^4)<\gamma_{\mathrm{bal}}(h_0)\le0\).
Hence \(\sigma^4\in[h_0,1]\), and root uniqueness implies
\[
h_\star(n,\sigma)=\sigma^4.
\]
The equality condition can be written as
\[
\log n+4\nu\log\sigma=\sigma^{-2}-1.
\]
Using \(\log(\sigma^{-2})=-2\log\sigma\) and
\(\sigma^2\sigma^{-2}=1\), it gives both scalar identities
\[
\begin{aligned}
\sigma^{-2}+\frac{3\nu}{2}\log(\sigma^{-2})
&=1+\log(n\sigma^\nu),\\
2\nu\log(\sigma^{-2})
+\sigma^{-2}\{1+\log(\sigma^2\sigma^{-2})\}
&=1+\log n.
\end{aligned}
\]
Thus \(z=\tau=\sigma^{-2}\), and both resolution formulas equal \(\sigma^4\).

At the direct--noisy equality \(\sigma=h_0\), the direct result gives
\[
h_\star(n,\sigma)=\sigma=h_0,\qquad
\log(n\sigma^\nu)=0.
\]
The intermediate scalar equation then holds at \(z=1\), and
\(\sigma z^{-3/2}=\sigma\).

Finally, \(c_{\mathrm{br}}\le1,c_\tau,c_{\mathrm{pow}}\) and
\(C_{\mathrm{br}}\ge A_{\mathrm{int}},C_\tau,C_{\mathrm{pow}}\).
All comparison factors above are positive. Replacing their lower constants
by \(c_{\mathrm{br}}\) and their upper constants by \(C_{\mathrm{br}}\)
therefore gives simultaneous branch comparisons, with the exact equations
and interfaces unchanged.
% lean: CausalSmith.Stat.RdTruesideNoiseFrontier.intermediateCondition_iff
% lean: CausalSmith.Stat.RdTruesideNoiseFrontier.intermediate_branch_representation
% lean: CausalSmith.Stat.RdTruesideNoiseFrontier.intermediate_branch_bounds
% lean: CausalSmith.Stat.RdTruesideNoiseFrontier.compact_branch_representation
% lean: CausalSmith.Stat.RdTruesideNoiseFrontier.compactScalar_bounds
% lean: CausalSmith.Stat.RdTruesideNoiseFrontier.compactPower_bounds
% lean: CausalSmith.Stat.RdTruesideNoiseFrontier.intermediate_compact_interface
% lean: CausalSmith.Stat.RdTruesideNoiseFrontier.direct_interface
% lean: CausalSmith.Stat.RdTruesideNoiseFrontier.frontier_branches

\item \textbf{Shrinking-noise sequences.}
Fix \(0<a<1/\nu\) and set \(\sigma=n^{-a}\). Define
\[
\eta_a=1-4a\nu,\qquad
\epsilon_a=\frac1{2(|\eta_a|+1)}.
\]
Since \(2a>0\), the relation \(\log n=o(n^{2a})\) and divergence of \(n^{2a}\)
give an integer \(N\ge2\) such that, for all \(n\ge N\),
\[
\log n\le\epsilon_a n^{2a},
\qquad n^{2a}\ge2.
\]
Consequently,
\[
\eta_a\log n
\le|\eta_a|\epsilon_a n^{2a}
\le\frac12n^{2a}
\le n^{2a}-1.
\]
The relevant quantities are exactly
\[
\log(n\sigma^{4\nu})=(1-4a\nu)\log n,
\qquad
\sigma^{-2}=n^{2a}.
\]
Thus the intermediate branch holds for \(n\ge N\); moreover,
\(0<\sigma\le1\) and \(\sigma>h_0\) for every \(n\ge2\).

Define
\[
d_{\mathrm{shrink}}=1-a\nu>0,\qquad
K_{\mathrm{shrink}}=d_{\mathrm{shrink}}+\frac1{\log2}>0.
\]
For \(n\ge2\),
\[
1+\log(n\sigma^\nu)=1+d_{\mathrm{shrink}}\log n,
\]
and \(\log n\ge\log2\) implies
\[
d_{\mathrm{shrink}}\log n
\le1+\log(n\sigma^\nu)
\le K_{\mathrm{shrink}}\log n.
\]
Applying the intermediate comparison with \(c_{\mathrm{br}},C_{\mathrm{br}}\)
therefore gives, for \(n\ge N\),
\[
\frac{c_{\mathrm{br}}}{K_{\mathrm{shrink}}^{3/2}}
\,n^{-a}(\log n)^{-3/2}
\le h_\star(n,n^{-a})
\le
\frac{C_{\mathrm{br}}}{d_{\mathrm{shrink}}^{3/2}}
\,n^{-a}(\log n)^{-3/2}.
\]
Taking the positive power \(\beta\) yields
\[
\begin{aligned}
\left(\frac{c_{\mathrm{br}}}{K_{\mathrm{shrink}}^{3/2}}\right)^\beta
n^{-a\beta}(\log n)^{-3\beta/2}
&\le r_\beta(n,n^{-a})\\
&\le
\left(\frac{C_{\mathrm{br}}}{d_{\mathrm{shrink}}^{3/2}}\right)^\beta
n^{-a\beta}(\log n)^{-3\beta/2}.
\end{aligned}
\]
Both constants are positive and may depend on \(a\).

For \(a\ge1/\nu\), monotonicity in the exponent for \(n>1\) gives
\[
0<n^{-a}\le n^{-1/\nu}=h_0\le1.
\]
Hence the exact direct branch applies for every \(n\ge2\), with
\[
h_\star(n,n^{-a})=h_0,
\qquad
r_\beta(n,n^{-a})=n^{-\beta/\nu}.
\]
% lean: CausalSmith.Stat.RdTruesideNoiseFrontier.eventually_shrinking_intermediate
% lean: CausalSmith.Stat.RdTruesideNoiseFrontier.shrinking_intermediate_rate_bounds
% lean: CausalSmith.Stat.RdTruesideNoiseFrontier.frontier_sequences

\item \textbf{Fixed positive noise.}
Fix \(0<\sigma\le1\), and define
\[
\begin{aligned}
\ell_{\mathrm{noise}}&=\log(\sigma^2),\\
M_{\mathrm{noise}}&=\max\left\{
1,\,-2\ell_{\mathrm{noise}},\,|\log(e+\sigma^2)|
\right\},\\
K_{\mathrm{noise}}&=1+|\log(e+\sigma^2)|,\\
T_{\mathrm{noise}}&=\sigma^{-2}-1-4\nu\log\sigma.
\end{aligned}
\]
Since \(\log n\to\infty\) and \(h_0\to0\), there is \(N\ge2\) such that, for
\(n\ge N\),
\[
e^{M_{\mathrm{noise}}}\le\log n,\qquad
T_{\mathrm{noise}}+1\le\log n,\qquad
h_0<\sigma.
\]
The middle inequality gives
\[
\log(n\sigma^{4\nu})
=\log n+4\nu\log\sigma
>\sigma^{-2}-1,
\]
so the compact branch applies.

Retain
\[
\Lambda_n=1+\log n,\qquad
\Omega_{n,\sigma}=\log(e+\sigma^2\Lambda_n),
\]
and define
\[
L_{\mathrm{log}}=\log\Lambda_n.
\]
The first threshold implies
\[
L_{\mathrm{log}}\ge M_{\mathrm{noise}}\ge1,
\qquad
L_{\mathrm{log}}\ge-2\ell_{\mathrm{noise}}.
\]
Monotonicity of the logarithm gives
\[
\Omega_{n,\sigma}
\ge\log(\sigma^2\Lambda_n)
=\ell_{\mathrm{noise}}+L_{\mathrm{log}}
\ge\frac12L_{\mathrm{log}}.
\]
Since \(\Lambda_n\ge1\),
\[
e+\sigma^2\Lambda_n\le(e+\sigma^2)\Lambda_n,
\]
and hence
\[
\begin{aligned}
\Omega_{n,\sigma}
&\le\log(e+\sigma^2)+L_{\mathrm{log}}\\
&\le|\log(e+\sigma^2)|L_{\mathrm{log}}+L_{\mathrm{log}}\\
&=K_{\mathrm{noise}}L_{\mathrm{log}}.
\end{aligned}
\]
Applying the compact rate comparison and taking the power \(2\beta>0\) gives
\[
\begin{aligned}
c_{\mathrm{br}}2^{-2\beta}
\left(\frac{L_{\mathrm{log}}}{\Lambda_n}\right)^{2\beta}
&\le r_\beta(n,\sigma)\\
&\le
C_{\mathrm{br}}K_{\mathrm{noise}}^{2\beta}
\left(\frac{L_{\mathrm{log}}}{\Lambda_n}\right)^{2\beta}.
\end{aligned}
\]
Since \(\Lambda_n=\log(en)\) and \(L_{\mathrm{log}}=\log\log(en)\), this is the
fixed-noise assertion with the permitted dependence on \(\sigma\).
% lean: CausalSmith.Stat.RdTruesideNoiseFrontier.fixed_noise_rate_bounds
% lean: CausalSmith.Stat.RdTruesideNoiseFrontier.frontier_sequences

\item \textbf{The case \(\beta=1\), \(\sigma=n^{-1/6}\).}
For every \(n\ge2\),
\[
\nu=3,\qquad h_0=n^{-1/3},\qquad
\sigma=n^{-1/6}>h_0,\qquad 0<\sigma\le1.
\]
Power identities give
\[
\sigma^{12}=n^{-2},\qquad
n\sigma^{12}=n^{-1},\qquad
\sigma^{-2}=n^{1/3}.
\]
Thus
\[
\log(n\sigma^{12})=-\log n
\le n^{1/3}-1=\sigma^{-2}-1,
\]
so the intermediate branch holds at every such sample size. Also
\[
n\sigma^3=n^{1/2},
\qquad
\log(n\sigma^3)=\frac12\log n,
\qquad
r_1(n,\sigma)=h_\star(n,\sigma).
\]
The intermediate comparison at \(\beta=1\) therefore yields
\[
r_1(n,n^{-1/6})
=h_\star(n,n^{-1/6})
\asymp
\frac{n^{-1/6}}{[1+\tfrac12\log n]^{3/2}},
\]
with constants independent of \(n\).
% lean: CausalSmith.Stat.RdTruesideNoiseFrontier.beta_one_frontier

\item \textbf{A simultaneous choice of constants.}
Finally, define
\[
c=\min\{c_{\mathrm{mat}},c_{\mathrm{br}}\},
\qquad
C=\max\{C_{\mathrm{mat}},C_{\mathrm{br}}\}.
\]
Then
\[
0<c,\qquad 0<C,\qquad
c\le c_{\mathrm{mat}},c_{\mathrm{br}},
\qquad
C\ge C_{\mathrm{mat}},C_{\mathrm{br}}.
\]
The root bounds give
\[
h_\star(n,\sigma)\ge h_0>0,
\qquad
r_\beta(n,\sigma)=h_\star(n,\sigma)^\beta\ge0.
\]
Thus, in the matched chains,
\[
c\,r_\beta(n,\sigma)
\le c_{\mathrm{mat}}r_\beta(n,\sigma),
\qquad
C_{\mathrm{mat}}r_\beta(n,\sigma)
\le C\,r_\beta(n,\sigma).
\]
The branch comparisons admit the same replacement: their factors
\[
\frac{\sigma}{[1+\log(n\sigma^\nu)]^{3/2}},
\qquad
\frac{\Lambda_n}{\Omega_{n,\sigma}},
\qquad
\left(\frac{\Omega_{n,\sigma}}{\Lambda_n}\right)^{2\beta}
\]
are positive on their respective noisy branches. Multiplication therefore
preserves the inequalities when their lower constants are reduced to \(c\)
and their upper constants are increased to \(C\). The root, branch equations,
uniqueness assertions, and interfaces retain their exact values. Together
with the sequence comparisons and the case \(\beta=1\), this proves every
assertion simultaneously.
% lean: CausalSmith.Stat.RdTruesideNoiseFrontier.matched_frontier_weaken
% lean: CausalSmith.Stat.RdTruesideNoiseFrontier.frontier_branches_weaken
% lean: CausalSmith.Stat.RdTruesideNoiseFrontier.uniform_frontier
\end{enumerate}
\end{proof}

\begin{proof}[Proof of \cref{obj:prop:published-class-converse-transfer}]
Fix \(\beta\in(0,1]\).

\begin{enumerate}
\item By \cref{obj:thm:uniform-frontier}, there is a constant \(c>0\), depending only on \(\beta\), such that, simultaneously for every \(n\ge2\) and \(\sigma\in[0,1]\),
\[
c\,r_\beta(n,\sigma)\le R_\beta(n,\sigma),
\qquad
c\,r_\beta(n,\sigma)\le\mathcal L_\beta(n,\sigma).
\]
We use this same constant for both transferred lower bounds.
% lean: CausalSmith.Stat.RdTruesideNoiseFrontier.uniform_frontier

\item Fix \(\sigma\in[0,1]\) and \(P\in\mathcal M_\beta(\sigma)\). Define the embedded law and its accompanying density and mean versions by
\[
Q=P,\qquad
(Z,V^0,V^1,E)=(X,Y^0,Y^1,\varepsilon),
\qquad
a(z)=f(P,z),\qquad v_d(z)=\mu_d(P,z)
\quad(d\in\{0,1\},\ z\in\mathbb R).
\]
The density identity is inherited from \cref{obj:synth_1}. For every measurable set \(F_{\mathrm{score}}\subseteq\mathbb R\), the conditional-mean identity becomes
\[
\begin{aligned}
\int_{\{Z\in F_{\mathrm{score}}\}}V^d\,dQ
&=\int_{F_{\mathrm{score}}\cap[-1,1]}
       \mu_d(P,z)f(P,z)\,dz\\
&=\int_{F_{\mathrm{score}}}\mu_d(P,z)f(P,z)\,dz\\
&=\int_{F_{\mathrm{score}}}v_d(z)a(z)\,dz.
\end{aligned}
\]
The first equality is supplied by \cref{obj:synth_1}; the second follows because \(f(P,z)=0\) outside \([-1,1]\). Thus the chosen \(v_d\) remain conditional-mean versions on the real score space.
% lean: CausalSmith.Stat.RdTruesideNoiseFrontier.benchmark_version_on_real

Choose
\[
r=\frac12.
\]
Since \((-r,r)\subseteq[-1,1]\), the full-interval continuity in \cref{obj:synth_1} gives continuity of \(a\) and both \(v_d\) on \((-r,r)\). Moreover, \cref{obj:ass:density-bounds,obj:ass:mean-bounds} give
\[
a(z)\ge\frac14>0\quad(z\in(-r,r)),
\qquad
0\le\frac14\le v_d(0)\le\frac34\le1
\quad(d\in\{0,1\}).
\]
The Gaussian condition follows from \cref{obj:ass:gaussian}:
\[
E\sim N(0,1).
\]
To verify the required joint independence, define the measurable map
\[
\Phi_{\mathrm{obs}}(z,y_0,y_1)
=
\left(
z,\,
\mathbf1\{z\ge0\}y_1+
\bigl(1-\mathbf1\{z\ge0\}\bigr)y_0
\right),
\qquad
(z,y_0,y_1)\in\mathbb R\times\{0,1\}^2.
\]
Then
\[
(Z,B)=\Phi_{\mathrm{obs}}(X,Y^0,Y^1).
\]
By \cref{obj:ass:error-independence}, independence is preserved under this measurable transformation, so
\[
E\perp(Z,B).
\]
All defining conditions of \(\mathcal G(\sigma)\) therefore hold.
% lean: CausalSmith.Stat.RdTruesideNoiseFrontier.benchmarkEmbedding_mem

The embedding leaves the coordinates used by the observation map unchanged. Consequently,
\[
\vartheta(Q)=v_1(0)-v_0(0)=\theta(P),
\qquad
Q\circ(R,A,B)^{-1}=P_{\mathrm{obs}},
\]
and hence, for every integer \(n\ge0\),
\[
\bigl(Q\circ(R,A,B)^{-1}\bigr)^{\otimes n}
\otimes\operatorname{Unif}[0,1]
=
P_{\mathrm{obs}}^{\otimes n}
\otimes\operatorname{Unif}[0,1]
=
\mathbb P_{P,n}.
\]
At \(n=0\), the sample factor is the empty product probability law, so the identity also covers that case.
% lean: CausalSmith.Stat.RdTruesideNoiseFrontier.benchmarkEmbedding_preserves

\item Fix \(n\ge2\) and \(\sigma\in[0,1]\). For every measurable estimate \(T\), target and experiment preservation give
\[
\sup_{P\in\mathcal M_\beta(\sigma)}
\mathbb E_{P,n}|T-\theta(P)|
\le
\sup_{Q\in\mathcal G(\sigma)}
\mathbb E_Q|T-\vartheta(Q)|.
\]
Taking infima over the common estimator class gives the corresponding comparison of extended nonnegative minimax objectives.

To justify the comparison after conversion to real values, note that every \(Q\in\mathcal G(\sigma)\) satisfies
\[
-1\le v_1(0)-v_0(0)=\vartheta(Q)\le1.
\]
Define the constant estimate by
\[
T_{\mathrm{zero}}(S_n,U)=0.
\]
Since the randomized experiments are probability laws,
\[
\sup_{Q\in\mathcal G(\sigma)}
\mathbb E_Q|T_{\mathrm{zero}}-\vartheta(Q)|
\le1.
\]
Thus the larger-class extended minimax risk is finite, and conversion to real values preserves the preceding inequality:
\[
R_\beta(n,\sigma)\le\mathcal R_{\mathcal G}(n,\sigma).
\]
% lean: CausalSmith.Stat.RdTruesideNoiseFrontier.benchmark_risk_le_gaussian

For interval length, define the constant interval procedure by
\[
I_{\mathrm{full}}(S_n,U)=[-1,1].
\]
The same target bounds imply, for every \(Q\in\mathcal G(\sigma)\),
\[
Q\{\vartheta(Q)\in I_{\mathrm{full}}\}=1,
\qquad
\mathbb E_Q\operatorname{length}(I_{\mathrm{full}})=2.
\]
Hence the larger-class honest interval class is nonempty and its extended minimax expected length is at most \(2\).

Now let \(I=[\ell,\upsilon]\) be any interval procedure honest on \(\mathcal G(\sigma)\). For every benchmark law \(P\), its embedded law \(Q\) satisfies
\[
\mathbb P_{P,n}\{\theta(P)\in I\}
=
Q\{\vartheta(Q)\in I\}
\ge0.9.
\]
Thus \(I\) belongs to the benchmark honest interval class of \cref{obj:def:honest-class}. Moreover,
\[
\mathcal L_\beta(n,\sigma)
\le
\sup_{P\in\mathcal M_\beta(\sigma)}
\mathbb E_{P,n}[\upsilon-\ell]
\le
\sup_{Q\in\mathcal G(\sigma)}
\mathbb E_Q[\upsilon-\ell].
\]
The second inequality follows from experiment preservation. Taking the infimum over intervals honest on \(\mathcal G(\sigma)\) proves the extended minimax comparison; the finite bound \(2\) permits its conversion to real values. Therefore
\[
\mathcal L_\beta(n,\sigma)\le\mathcal L_{\mathcal G}(n,\sigma).
\]
Combining these comparisons with the lower bounds already obtained yields
\[
c\,r_\beta(n,\sigma)
\le R_\beta(n,\sigma)
\le\mathcal R_{\mathcal G}(n,\sigma),
\qquad
c\,r_\beta(n,\sigma)
\le\mathcal L_\beta(n,\sigma)
\le\mathcal L_{\mathcal G}(n,\sigma).
\]
% lean: CausalSmith.Stat.RdTruesideNoiseFrontier.benchmark_length_le_gaussian

\item Finally, fix \(\sigma\in(0,1]\) and \(Q\in\mathcal G(\sigma)\). Its joint independence condition, followed by the measurable scaling map \(e\mapsto\sigma e\), gives
\[
E\perp(Z,B),
\qquad
\sigma E\perp(Z,B).
\]
These supply the standardized and scaled error independence in the stated Assumption 1Y condition. The local density condition supplies a radius \(r>0\) for which
\[
a(z)>0\qquad(z\in(-r,r)),
\]
which is Assumption 2C. Finally, scaling the standard Gaussian error gives
\[
N(0,1)\circ(e\mapsto\sigma e)^{-1}=N(0,\sigma^2),
\qquad
\sigma E\sim N(0,\sigma^2).
\]
This is Assumption 7, completing the three asserted conditions.
% lean: CausalSmith.Stat.RdTruesideNoiseFrontier.gaussian_published_hypotheses
\end{enumerate}
\end{proof}

\begin{proof}[Proof of \cref{obj:lem:classical-legendre-facts}]
Write the shifted Legendre polynomials in the normalization
\[
R_j(x):=\sum_{m=0}^{j}(-1)^m
  \binom{j}{m}\binom{j+m}{m}x^m,
\qquad
P_j^{\mathrm{Leg}}(t)=R_j\!\left(\frac{1-t}{2}\right),
\qquad j\in\mathbb N,\quad x,t\in\mathbb R.
\]

\begin{enumerate}
\item \emph{Orthogonality.}
The affine substitution \(t=1-2x\), including reversal of the integration endpoints, gives
\[
\int_0^1R_j(x)R_k(x)\,dx
=-\frac12\int_1^{-1}
  P_j^{\mathrm{Leg}}(t)P_k^{\mathrm{Leg}}(t)\,dt
=\frac12\int_{-1}^1
  P_j^{\mathrm{Leg}}(t)P_k^{\mathrm{Leg}}(t)\,dt.
\]
We therefore compute the inner product on the unit interval.

For \(k\in\mathbb N\), define
\[
\mathcal B_k(x):=x^k(1-x)^k,
\qquad x\in\mathbb R.
\]
Expanding \((1-x)^k\) and differentiating \(k\) times gives Rodrigues' identity:
\[
\mathcal B_k^{(k)}(x)
=\sum_{m=0}^k(-1)^m\binom{k}{m}
  \frac{(k+m)!}{m!}x^m
=k!R_k(x).
\]
For every \(0\le i<k\), the factors \(x^k\) and \((1-x)^k\), respectively, ensure that
\[
\mathcal B_k^{(i)}(0)=\mathcal B_k^{(i)}(1)=0.
\]
Indeed, in each term of the Leibniz expansion, fewer than \(k\) derivatives fall on the factor vanishing at the relevant endpoint. Consequently, repeated integration by parts gives, for every polynomial \(R_j\),
\[
\int_0^1R_j(x)\mathcal B_k^{(k)}(x)\,dx
=(-1)^k\int_0^1R_j^{(k)}(x)\mathcal B_k(x)\,dx.
\]
For \(k=0\), this identity involves zero integrations by parts and has no boundary conditions.
% lean: Causalean.Mathlib.Analysis.SpecialFunctions.Jacobi.integral_mul_iterated_deriv_eq

First consider \(j=k\). The coefficient of \(x^k\) in \(R_k\) gives
\[
R_k^{(k)}(x)=k!(-1)^k\binom{2k}{k}.
\]
Combining this identity with Rodrigues' formula and integration by parts yields
\[
\begin{aligned}
k!\int_0^1R_k(x)^2\,dx
&=\int_0^1R_k(x)\mathcal B_k^{(k)}(x)\,dx\\
&=(-1)^k\int_0^1R_k^{(k)}(x)\mathcal B_k(x)\,dx\\
&=k!\binom{2k}{k}\int_0^1x^k(1-x)^k\,dx.
\end{aligned}
\]
The beta integral evaluates to
\[
\int_0^1x^k(1-x)^k\,dx
=\frac{k!}{(k+1)\prod_{i=0}^{k-1}(k+2+i)}
=\frac{(k!)^2}{(2k+1)!},
\]
where the product is \(1\) when \(k=0\). Since \(k!>0\), cancellation and the factorial formula for the binomial coefficient give
\[
\int_0^1R_k(x)^2\,dx
=\binom{2k}{k}\frac{(k!)^2}{(2k+1)!}
=\frac{1}{2k+1}.
\]
% lean: Causalean.Mathlib.Analysis.SpecialFunctions.Jacobi.shiftedLegendre_sq_integral_unit

If \(j<k\), the degree of \(R_j\) is \(j\), so \(R_j^{(k)}=0\). Thus
\[
k!\int_0^1R_j(x)R_k(x)\,dx
=(-1)^k\int_0^1R_j^{(k)}(x)\mathcal B_k(x)\,dx
=0,
\]
and the inner product vanishes. If \(k<j\), interchange the two factors and apply the same calculation. Hence
\[
\int_0^1R_j(x)R_k(x)\,dx
=
\begin{cases}
\dfrac{1}{2j+1},&j=k,\\
0,&j\ne k.
\end{cases}
\]
Multiplying the affine-substitution identity by \(2\) proves the asserted inner product on \([-1,1]\).
% lean: CausalSmith.Stat.RdTruesideNoiseFrontier.legendreP_orthogonal_exact

\item \emph{Reflection symmetry.}
The shifted Legendre reflection identity is
\[
R_j(x)=(-1)^jR_j(1-x),
\qquad x\in\mathbb R.
\]
Applying it at \(x=(1+t)/2\) gives
\[
P_j^{\mathrm{Leg}}(-t)
=R_j\!\left(\frac{1+t}{2}\right)
=(-1)^jR_j\!\left(\frac{1-t}{2}\right)
=(-1)^jP_j^{\mathrm{Leg}}(t),
\qquad t\in\mathbb R.
\]
% lean: CausalSmith.Stat.RdTruesideNoiseFrontier.legendreP_reflection_exact

\item \emph{The bound on the interval.}
The finite sum defining \(R_j\) is also the normalized shifted Jacobi polynomial with both shape parameters equal to zero: its coefficient identity is
\[
(-1)^m\binom jm
\frac{(j+m)!/j!}{m!}
=(-1)^m\binom jm\binom{j+m}{m},
\qquad 0\le m\le j.
\]
For \(\xi\in\mathbb R\) and \(r\in\mathbb N\), define the rising and falling factorial products by
\[
(\xi)_{\uparrow r}:=\prod_{i=0}^{r-1}(\xi+i),
\qquad
(\xi)_{\downarrow r}:=\prod_{i=0}^{r-1}(\xi-i),
\]
with both empty products equal to \(1\). Their definitions give
\[
(-\xi)_{\downarrow r}=(-1)^r(\xi)_{\uparrow r},
\qquad
(\xi)_{\uparrow m}(\xi+m)_{\uparrow(j-m)}
=(\xi)_{\uparrow j},
\quad 0\le m\le j.
\]
The falling-factorial convolution is
\[
(\xi+\eta)_{\downarrow j}
=\sum_{m=0}^j\binom jm
  (\xi)_{\downarrow m}(\eta)_{\downarrow(j-m)},
\qquad \xi,\eta\in\mathbb R.
\]
For nonnegative integer arguments, comparing coefficients of degree \(j\) in
\[
(1+\zeta)^{\xi+\eta}
=(1+\zeta)^\xi(1+\zeta)^\eta,
\qquad \xi,\eta\in\mathbb N,\quad \zeta\in\mathbb R,
\]
and multiplying by \(j!\) proves the convolution. Both sides are polynomials in the two arguments, so the identity extends to all real arguments.

Specialize the convolution to \(\xi=-(j+1)\) and \(\eta=j\). The product identities yield, for \(0\le m\le j\),
\[
\begin{aligned}
(-1)_{\downarrow j}&=(-1)^j j!,\\
(-(j+1))_{\downarrow m}
&=(-1)^m(j+1)_{\uparrow m}
=(-1)^m\frac{(j+m)!}{j!},\\
(j)_{\downarrow(j-m)}
&=(m+1)_{\uparrow(j-m)}
=\frac{j!}{m!}.
\end{aligned}
\]
Consequently the convolution becomes
\[
(-1)^j j!
=\sum_{m=0}^j\binom jm(-1)^m
  \frac{(j+m)!}{j!}\frac{j!}{m!}
=\sum_{m=0}^j(-1)^m\binom jm\frac{(j+m)!}{m!}.
\]
Dividing by \(j!>0\) supplies the normalization:
\[
R_j(1)
=\sum_{m=0}^j(-1)^m\binom jm\binom{j+m}{m}
=\frac{1}{j!}\sum_{m=0}^j(-1)^m\binom jm\frac{(j+m)!}{m!}
=(-1)^j,
\qquad |R_j(1)|=1.
\]
These identities include \(j=0\), when the products are empty and \(0!=1\).
At the other endpoint, direct evaluation of the defining sum gives
\[
R_j(0)=1,
\qquad |R_j(0)|=1.
\]
% lean: Causalean.Mathlib.Analysis.SpecialFunctions.Jacobi.jacobiShifted_one

When \(j=0\), the defining sum gives \(R_0(x)=1\) for every real \(x\), proving the required bound. For \(j\ge1\), define
\[
\lambda_j:=j(j+1)>0,
\qquad
\mathcal E_j(x):=
R_j(x)^2+\frac{x(1-x)R_j'(x)^2}{\lambda_j},
\qquad x\in\mathbb R.
\]
We derive the differential identity governing this energy. Define the coefficients over their full range by
\[
c_{j,m}:=
\begin{cases}
(-1)^m\binom jm\binom{j+m}{m},&0\le m\le j,\\
0,&m>j,
\end{cases}
\qquad m\in\mathbb N.
\]
The binomial coefficient formulas imply
\[
(m+1)^2c_{j,m+1}
=-(j-m)(j+m+1)c_{j,m},
\qquad 0\le m<j.
\]
The coefficient of \(x^m\) in the polynomial
\[
x(1-x)R_j''(x)+(1-2x)R_j'(x)+\lambda_jR_j(x)
\]
is
\[
(m+1)^2c_{j,m+1}
+\bigl(\lambda_j-m(m+1)\bigr)c_{j,m}.
\]
It vanishes for \(m<j\) by the recurrence, for \(m=j\) because \(c_{j,j+1}=0\) and \(\lambda_j=j(j+1)\), and for \(m>j\) because both coefficients are zero. Therefore
\[
x(1-x)R_j''(x)+(1-2x)R_j'(x)+\lambda_jR_j(x)=0,
\qquad x\in\mathbb R.
\]
% lean: Causalean.Mathlib.Analysis.SpecialFunctions.Jacobi.jacobiShifted_ode

Differentiation of the energy, followed by substitution of this differential equation, yields
\[
\begin{aligned}
\mathcal E_j'(x)
&=2R_j(x)R_j'(x)
+\frac{(1-2x)R_j'(x)^2
       +2x(1-x)R_j'(x)R_j''(x)}{\lambda_j}\\
&=(2x-1)\frac{R_j'(x)^2}{\lambda_j}.
\end{aligned}
\]
Thus the energy is nonincreasing on \([0,1/2]\) and nondecreasing on \([1/2,1]\). Its endpoint values are
\[
\mathcal E_j(0)=R_j(0)^2=1,
\qquad
\mathcal E_j(1)=R_j(1)^2=1.
\]
For \(x\le1/2\), compare with the left endpoint; for \(x>1/2\), compare with the right endpoint. These comparisons give
\[
\mathcal E_j(x)\le
\begin{cases}
\mathcal E_j(0),&0\le x\le1/2,\\
\mathcal E_j(1),&1/2<x\le1,
\end{cases}
\qquad
\mathcal E_j(x)\le1.
\]
Since the additional term in the energy is nonnegative on the unit interval,
\[
R_j(x)^2
\le R_j(x)^2+\frac{x(1-x)R_j'(x)^2}{\lambda_j}
=\mathcal E_j(x)\le1,
\qquad 0\le x\le1.
\]
Hence \(|R_j(x)|\le1\) throughout that interval. Finally,
\[
-1\le t\le1
\quad\Longrightarrow\quad
0\le\frac{1-t}{2}\le1
\quad\Longrightarrow\quad
|P_j^{\mathrm{Leg}}(t)|
=\left|R_j\!\left(\frac{1-t}{2}\right)\right|\le1.
\]
% lean: CausalSmith.Stat.RdTruesideNoiseFrontier.legendreP_abs_le_one_exact

\item \emph{The endpoint value.}
Evaluation at the positive endpoint gives
\[
P_j^{\mathrm{Leg}}(1)=R_j(0)=1.
\]
Applying reflection symmetry at \(t=1\) therefore yields
\[
P_j^{\mathrm{Leg}}(-1)
=(-1)^jP_j^{\mathrm{Leg}}(1)
=(-1)^j.
\]
% lean: CausalSmith.Stat.RdTruesideNoiseFrontier.legendreP_neg_one_exact
\end{enumerate}
\end{proof}

\begin{proof}[Proof of \cref{obj:thm:direct-reduction}]
Fix \(0<\beta\le1\). For integers \(n\ge2\), define
\[
r_{\mathrm{dir},n}:=n^{-\beta/(2\beta+1)}.
\]
Expectations of nonnegative losses and their suprema below are taken in the extended nonnegative reals.

\begin{enumerate}
\item \textbf{The selected radius at zero noise.}
By \cref{obj:lem:positive-endpoint}, there is a constant \(K_{\mathrm{end}}>0\), depending only on \(\beta\), such that, for every integer \(L\ge1\),
\[
M_{L,\beta}\le K_{\mathrm{end}}L^{-2\beta},
\qquad
\int_0^1 K_L(x)^2\,dx\le K_{\mathrm{end}}L^2.
\]
Define
\[
K_{\mathrm{cert}}
:=48K_{\mathrm{end}}+28\sqrt{30K_{\mathrm{end}}}>0.
\]
For a fixed \(n\ge2\), put
\[
a_{\mathrm{deg}}:=n^{1/(4\beta+2)},
\qquad
L_{\mathrm{dir}}:=\lfloor a_{\mathrm{deg}}\rfloor,
\qquad
L_{\max}:=\lceil a_{\mathrm{deg}}\rceil.
\]
Since \(a_{\mathrm{deg}}\ge1\), the floor inequalities give
\[
1\le L_{\mathrm{dir}},
\qquad
a_{\mathrm{deg}}<L_{\mathrm{dir}}+1\le2L_{\mathrm{dir}},
\qquad
L_{\mathrm{dir}}\le a_{\mathrm{deg}},
\qquad
L_{\mathrm{dir}}\le L_{\max}.
\]
These inequalities apply throughout \(a_{\mathrm{deg}}\ge1\), including \(1\le a_{\mathrm{deg}}<2\). Consequently,
\[
\begin{aligned}
M_{L_{\mathrm{dir}},\beta}
&\le K_{\mathrm{end}}L_{\mathrm{dir}}^{-2\beta}\\
&\le K_{\mathrm{end}}(a_{\mathrm{deg}}/2)^{-2\beta}\\
&=K_{\mathrm{end}}a_{\mathrm{deg}}^{-2\beta}2^{2\beta}
\le4K_{\mathrm{end}}a_{\mathrm{deg}}^{-2\beta},
\end{aligned}
\]
where the last inequality uses \(\beta\le1\).

In the variance formula of \cref{obj:synth_6}, the coefficient of the zeroth derivative is one at \(\sigma=0\), and every positive-order coefficient vanishes. Thus
\[
V_L(0)=\frac34\int_0^1K_L(x)^2\,dx,
\]
and hence
\[
\frac{40V_{L_{\mathrm{dir}}}(0)}{n}
\le\frac{30K_{\mathrm{end}}L_{\mathrm{dir}}^2}{n}
\le\frac{30K_{\mathrm{end}}a_{\mathrm{deg}}^2}{n}.
\]
Taking square roots yields
\[
\sqrt{\frac{40V_{L_{\mathrm{dir}}}(0)}{n}}
\le\sqrt{30K_{\mathrm{end}}}\,
\frac{a_{\mathrm{deg}}}{\sqrt n}.
\]
The two balancing identities are
\[
a_{\mathrm{deg}}^{-2\beta}
=n^{-2\beta/(4\beta+2)}
=r_{\mathrm{dir},n},
\qquad
\frac{a_{\mathrm{deg}}}{\sqrt n}
=n^{1/(4\beta+2)-1/2}
=r_{\mathrm{dir},n}.
\]
The positive-index radius in \cref{obj:synth_6} therefore satisfies
\[
\rho_{L_{\mathrm{dir}}}
=12M_{L_{\mathrm{dir}},\beta}
+28\sqrt{\frac{40V_{L_{\mathrm{dir}}}(0)}{n}}
\le K_{\mathrm{cert}}r_{\mathrm{dir},n}.
\]
Since \(L_{\mathrm{dir}}\) belongs to the candidate set, the minimization in \cref{obj:def:public-selection} gives
\[
E_\beta(n,0)=\rho_{L_\star}
\le\rho_{L_{\mathrm{dir}}}
\le K_{\mathrm{cert}}r_{\mathrm{dir},n}.
\]
% lean: CausalSmith.Stat.RdTruesideNoiseFrontier.certificate_zero_rate_bound

\item \textbf{Direct witnesses and their resolution.}
Fix \(n\ge2\) and \(\sigma\in[0,1]\), and define
\[
h_0:=n^{-1/(2\beta+1)}.
\]
Because \(n\ge1\) and \(2\beta+1>0\),
\[
0<h_0\le1,
\qquad
nh_0^{2\beta+1}=1,
\qquad
h_0^\beta=r_{\mathrm{dir},n}.
\]

For \(0<h\le1\), use the laws \(P_h^{\mathrm{dir},\pm}\) of \cref{obj:def:direct-alternatives}. Define their mean perturbation by
\[
\delta_h(x):=\kappa h^\beta v_h^{\mathrm{dir}}(x),
\qquad x\in[-1,1],
\qquad
\kappa=\frac1{100}.
\]
The profile bounds in \cref{obj:synth_7} imply
\[
0\le\delta_h(x)\le\kappa,
\qquad
|\delta_h(x)-\delta_h(t)|
\le\kappa|x-t|^\beta
\quad(x,t\in[-1,1]).
\]
Together with the conditional means in \cref{obj:def:direct-alternatives}, these give
\[
\frac14
\le\mu_d(P_h^{\mathrm{dir},\pm},x)
\le\frac34,
\qquad
[\mu_0(P_h^{\mathrm{dir},\pm},\cdot)]_\beta=0,
\qquad
[\mu_1(P_h^{\mathrm{dir},\pm},\cdot)]_\beta\le\kappa\le2.
\]
Their score density is continuous and satisfies
\[
f(P_h^{\mathrm{dir},\pm},x)=\frac12
\quad(x\in[-1,1]),
\qquad
\int_{-1}^1f(P_h^{\mathrm{dir},\pm},x)\,dx=1.
\]
Their Gaussian error and independence properties are part of their construction. Thus all the conditions of \cref{obj:def:model} hold:
\[
P_h^{\mathrm{dir},\pm}\in\mathcal M_\beta(\sigma)
\quad(0<h\le1,\ 0\le\sigma\le1).
\]
Since \(v_h^{\mathrm{dir}}(0)=1\), their targets satisfy
\[
\theta(P_h^{\mathrm{dir},+})=\kappa h^\beta,
\qquad
\theta(P_h^{\mathrm{dir},-})=-\kappa h^\beta,
\qquad
\left|\theta(P_h^{\mathrm{dir},+})
-\theta(P_h^{\mathrm{dir},-})\right|
=2\kappa h^\beta.
\]
% lean: CausalSmith.Stat.RdTruesideNoiseFrontier.directAltLaw_model
% lean: CausalSmith.Stat.RdTruesideNoiseFrontier.directResolution_balance
% lean: CausalSmith.Stat.RdTruesideNoiseFrontier.directAltLaw_separation

\item \textbf{Noiseless information bound.}
For \(0<h\le1\), define the one-observation laws
\[
\mathsf P_{\sigma,h}^{\pm}
:=\operatorname{Law}_{P_h^{\mathrm{dir},\pm}}
(X+\sigma\varepsilon,D,Y).
\]
For probability measures \(\mathsf Q_1,\mathsf Q_2\) on a common measurable space, use
\[
\operatorname{TV}(\mathsf Q_1,\mathsf Q_2)
:=\sup_{\mathsf F\ \mathrm{measurable}}
|\mathsf Q_1(\mathsf F)-\mathsf Q_2(\mathsf F)|.
\]
In the absolutely continuous setting used here, the chi-squared divergence is
\[
\chi^2(\mathsf Q_1,\mathsf Q_2)
:=\int
\left(\frac{d\mathsf Q_1}{d\mathsf Q_2}-1\right)^2
\,d\mathsf Q_2.
\]

Let the reference measure on \(\mathbb R\times\{0,1\}^2\) be defined, for every measurable \(\mathsf F\), by
\[
\eta_{\mathrm{obs}}(\mathsf F)
:=\sum_{d=0}^1\sum_{y=0}^1
\int_{\mathbb R}\mathbf1_{\mathsf F}(x,d,y)\,dx.
\]
The densities
\[
p_h^\pm:=\frac{d\mathsf P_{0,h}^\pm}{d\eta_{\mathrm{obs}}}
\]
are given by
\[
\begin{aligned}
p_h^\pm(x,1,1)&=\frac14\pm\frac{\delta_h(x)}2,
&&0\le x\le1,\\
p_h^\pm(x,1,0)&=\frac14\mp\frac{\delta_h(x)}2,
&&0\le x\le1,\\
p_h^\pm(x,0,y)&=\frac14,
&&-1\le x<0,\quad y\in\{0,1\},
\end{aligned}
\]
and are zero elsewhere. Define a likelihood ratio on the entire observation space by
\[
\Lambda_h(x,d,y):=
\begin{cases}
p_h^+(x,1,y)/p_h^-(x,1,y),
&d=1,\ 0\le x\le1,\\
1,&\text{otherwise}.
\end{cases}
\]
The treated-cell densities lie between \(1/8\) and \(3/8\), so
\[
d\mathsf P_{0,h}^+
=\Lambda_h\,d\mathsf P_{0,h}^-,
\qquad
0\le\Lambda_h\le3,
\qquad
(\Lambda_h-1)^2\le4.
\]
In particular,
\[
\mathsf P_{0,h}^+\ll\mathsf P_{0,h}^-,
\qquad
\int(\Lambda_h-1)^2\,d\mathsf P_{0,h}^-<\infty.
\]

For each treated cell,
\[
p_h^+(x,1,y)-p_h^-(x,1,y)
=(2y-1)\delta_h(x),
\qquad
p_h^-(x,1,y)\ge\frac18.
\]
Consequently,
\[
0\le
p_h^-(x,1,y)(\Lambda_h(x,1,y)-1)^2
=\frac{\delta_h(x)^2}{p_h^-(x,1,y)}
\le8\delta_h(x)^2.
\]
Outside the treated cells, the corresponding integrand is zero. Summing over the two outcome cells gives
\[
\chi^2(\mathsf P_{0,h}^+,\mathsf P_{0,h}^-)
\le16\kappa^2h^{2\beta}
\int_0^1[v_h^{\mathrm{dir}}(x)]^2\,dx.
\]
The profile is \(1-x/h\) on \([0,h]\) and zero on \([h,1]\). With the substitution \(t=1-x/h\),
\[
\int_0^1[v_h^{\mathrm{dir}}(x)]^2\,dx
=\int_0^h(1-x/h)^2\,dx
=h\int_0^1t^2\,dt
=\frac h3.
\]
It follows that
\[
\chi^2(\mathsf P_{0,h}^+,\mathsf P_{0,h}^-)
\le\frac{16}{3}\kappa^2h^{2\beta+1}.
\]
At \(h=h_0\), the balance identity therefore yields
\[
n\chi^2(\mathsf P_{0,h_0}^+,\mathsf P_{0,h_0}^-)
\le\frac{16}{3}\kappa^2
=\frac{16}{30000}
\le\frac1{100}.
\]
% lean: CausalSmith.Stat.RdTruesideNoiseFrontier.directAltLaw_observed_ac_zero
% lean: CausalSmith.Stat.RdTruesideNoiseFrontier.directAltLaw_observed_likelihood_square_integrable_zero
% lean: CausalSmith.Stat.RdTruesideNoiseFrontier.directAltLaw_observed_chiSqDiv_zero_le

\item \textbf{Product testing bound and Gaussian contraction.}
Define the sample laws by
\[
\mathsf S_{\sigma,h}^{\pm}
:=(\mathsf P_{\sigma,h}^{\pm})^{\otimes n},
\]
and abbreviate the one-observation divergence at the selected resolution as
\[
\chi_{\mathrm{dir}}
:=\chi^2(\mathsf P_{0,h_0}^+,\mathsf P_{0,h_0}^-).
\]
The product likelihood ratio is
\[
\frac{d\mathsf S_{0,h_0}^+}{d\mathsf S_{0,h_0}^-}
\bigl((o_i)_{i=1}^n\bigr)
=\prod_{i=1}^n\Lambda_{h_0}(o_i),
\qquad
(o_i)_{i=1}^n\in(\mathbb R\times\{0,1\}^2)^n,
\]
almost everywhere under \(\mathsf S_{0,h_0}^-\). Since the one-observation likelihood ratio integrates to one,
\[
\int\Lambda_{h_0}\,d\mathsf P_{0,h_0}^-=1,
\qquad
\chi_{\mathrm{dir}}
=\int(\Lambda_{h_0}-1)^2\,d\mathsf P_{0,h_0}^-
=\int\Lambda_{h_0}^2\,d\mathsf P_{0,h_0}^- -1.
\]
Product integration likewise gives
\[
\int\prod_{i=1}^n\Lambda_{h_0}(o_i)
\,d\mathsf S_{0,h_0}^-\bigl((o_i)_{i=1}^n\bigr)
=\left(\int\Lambda_{h_0}\,d\mathsf P_{0,h_0}^-\right)^n
=1.
\]
Expanding the squared likelihood deviation and factoring its second moment over the product measure therefore yields
\[
\begin{aligned}
1+\chi^2(\mathsf S_{0,h_0}^+,\mathsf S_{0,h_0}^-)
&=\int\left(\prod_{i=1}^n\Lambda_{h_0}(o_i)\right)^2
\,d\mathsf S_{0,h_0}^-\bigl((o_i)_{i=1}^n\bigr)\\
&=\left(\int\Lambda_{h_0}^2\,d\mathsf P_{0,h_0}^-\right)^n\\
&=(1+\chi_{\mathrm{dir}})^n.
\end{aligned}
\]
Since \(\chi_{\mathrm{dir}}\ge0\) and \(n\chi_{\mathrm{dir}}\le1/100\),
\[
(1+\chi_{\mathrm{dir}})^n
\le e^{n\chi_{\mathrm{dir}}}
\le e^{1/100}
\le\sum_{j=0}^\infty(1/100)^j
=\frac{100}{99}.
\]
Here the penultimate inequality follows from the exponential series and \(j!\ge1\). Thus
\[
\chi^2(\mathsf S_{0,h_0}^+,\mathsf S_{0,h_0}^-)
\le\frac1{99}\le\frac4{25}.
\]
For absolutely continuous probability measures, the signed density difference has integral zero, and Cauchy--Schwarz gives
\[
\operatorname{TV}(\mathsf Q_1,\mathsf Q_2)
=\frac12\int
\left|\frac{d\mathsf Q_1}{d\mathsf Q_2}-1\right|
\,d\mathsf Q_2
\le\frac12\sqrt{\chi^2(\mathsf Q_1,\mathsf Q_2)}.
\]
Applying this to the product laws gives
\[
\operatorname{TV}(\mathsf S_{0,h_0}^+,\mathsf S_{0,h_0}^-)
\le\frac12\sqrt{\frac4{25}}=\frac15.
\]

For a sample
\[
\mathbf s=((x_i,d_i,y_i))_{i=1}^n
\in(\mathbb R\times\{0,1\}^2)^n,
\]
define the common Gaussian observation kernel by
\[
\mathsf K_{\sigma,n}(\mathbf s,\mathsf F)
:=
\Pr\!\left\{
((x_i+\sigma\varepsilon_i,d_i,y_i))_{i=1}^n
\in\mathsf F
\right\},
\qquad
(\varepsilon_i)_{i=1}^n\sim N(0,1)^{\otimes n}.
\]
The Gaussian variables in this definition are independent of the input sample. Gaussian error independence under the direct laws implies
\[
\mathsf S_{\sigma,h_0}^{\pm}
=\mathsf K_{\sigma,n}\mathsf S_{0,h_0}^{\pm}.
\]
For every measurable output event \(\mathsf F\), the function
\(\mathsf K_{\sigma,n}(\cdot,\mathsf F)\) takes values in \([0,1]\). The layer-cake formula therefore bounds the difference of its integrals by
\[
\begin{aligned}
&\left|
\int\mathsf K_{\sigma,n}(\mathbf s,\mathsf F)
\,d\mathsf S_{0,h_0}^+(\mathbf s)
-
\int\mathsf K_{\sigma,n}(\mathbf s,\mathsf F)
\,d\mathsf S_{0,h_0}^-(\mathbf s)
\right|\\
&\qquad\le
\int_0^1
\left|
\mathsf S_{0,h_0}^+
\{\mathsf K_{\sigma,n}(\cdot,\mathsf F)>t\}
-
\mathsf S_{0,h_0}^-
\{\mathsf K_{\sigma,n}(\cdot,\mathsf F)>t\}
\right|\,dt\\
&\qquad\le
\operatorname{TV}(\mathsf S_{0,h_0}^+,\mathsf S_{0,h_0}^-).
\end{aligned}
\]
Taking the supremum over output events proves
\[
\operatorname{TV}(\mathsf S_{\sigma,h_0}^+,\mathsf S_{\sigma,h_0}^-)
\le
\operatorname{TV}(\mathsf S_{0,h_0}^+,\mathsf S_{0,h_0}^-)
\le\frac15.
\]
The construction and this inequality also apply at \(\sigma=0\).
% lean: Causalean.Stat.iid_tv_le_one_fifth_of_chiSqDiv
% lean: CausalSmith.Stat.RdTruesideNoiseFrontier.sampleLaw_tv_le_zero_of_model

\item \textbf{Transfer to estimation risk and honest length.}
Let
\[
\mathsf U_{\mathrm{seed}}:=\operatorname{Uniform}[0,1],
\qquad
\mathsf E_{\sigma,h_0}^{\pm}
:=\mathsf S_{\sigma,h_0}^{\pm}\otimes\mathsf U_{\mathrm{seed}}.
\]
The product total-variation inequality gives
\[
\begin{aligned}
\operatorname{TV}(\mathsf E_{\sigma,h_0}^+,\mathsf E_{\sigma,h_0}^-)
&\le
\operatorname{TV}(\mathsf S_{\sigma,h_0}^+,\mathsf S_{\sigma,h_0}^-)
+\operatorname{TV}(\mathsf U_{\mathrm{seed}},\mathsf U_{\mathrm{seed}})\\
&\le\frac15.
\end{aligned}
\]
Define the common submeasure by
\[
d\mathsf H_{\sigma,h_0}
:=
\min\!\left\{
\frac{d\mathsf E_{\sigma,h_0}^+}
{d(\mathsf E_{\sigma,h_0}^++\mathsf E_{\sigma,h_0}^-)},
\frac{d\mathsf E_{\sigma,h_0}^-}
{d(\mathsf E_{\sigma,h_0}^++\mathsf E_{\sigma,h_0}^-)}
\right\}
d(\mathsf E_{\sigma,h_0}^++\mathsf E_{\sigma,h_0}^-).
\]
It satisfies
\[
\mathsf H_{\sigma,h_0}\le\mathsf E_{\sigma,h_0}^{\pm},
\qquad
\int d\mathsf H_{\sigma,h_0}
=
1-\operatorname{TV}(\mathsf E_{\sigma,h_0}^+,\mathsf E_{\sigma,h_0}^-)
\ge\frac45.
\]
The mass identity follows from
\(\min\{u,v\}=(u+v-|u-v|)/2\).

For any admissible estimator \(\widehat\theta\), domination by the common submeasure and the triangle inequality imply
\[
\begin{aligned}
&\mathbb E_{P_{h_0}^{\mathrm{dir},+},n}
|\widehat\theta-\theta(P_{h_0}^{\mathrm{dir},+})|
+
\mathbb E_{P_{h_0}^{\mathrm{dir},-},n}
|\widehat\theta-\theta(P_{h_0}^{\mathrm{dir},-})|\\
&\quad\ge
\int\left(
|\widehat\theta-\theta(P_{h_0}^{\mathrm{dir},+})|
+
|\widehat\theta-\theta(P_{h_0}^{\mathrm{dir},-})|
\right)\,d\mathsf H_{\sigma,h_0}\\
&\quad\ge
2\kappa h_0^\beta\int d\mathsf H_{\sigma,h_0}
\ge\frac85\kappa h_0^\beta.
\end{aligned}
\]
At least one of the two extended risks is therefore at least
\(\frac45\kappa h_0^\beta\), and both laws belong to the model class. Hence
\[
\inf_{\widehat\theta}
\sup_{P\in\mathcal M_\beta(\sigma)}
\mathbb E_{P,n}|\widehat\theta-\theta(P)|
\ge\frac45\kappa h_0^\beta.
\]

For any \(I\in\mathcal I_{\beta,n,\sigma}\), define its two coverage events by
\[
\mathsf F_{h_0,I}^{\pm}
:=
\{(S_n,U):\theta(P_{h_0}^{\mathrm{dir},\pm})\in I(S_n,U)\}.
\]
Honesty as defined in \cref{obj:def:honest-class} gives
\[
\mathsf E_{\sigma,h_0}^{+}(\mathsf F_{h_0,I}^{+})\ge\frac9{10},
\qquad
\mathsf E_{\sigma,h_0}^{-}(\mathsf F_{h_0,I}^{-})\ge\frac9{10}.
\]
Total variation transfers the first probability to the minus law:
\[
\mathsf E_{\sigma,h_0}^{-}(\mathsf F_{h_0,I}^{+})
\ge\frac9{10}-\frac15=\frac7{10}.
\]
Using the union bound under that law,
\[
\mathsf E_{\sigma,h_0}^{-}
(\mathsf F_{h_0,I}^{+}\cap\mathsf F_{h_0,I}^{-})
\ge\frac7{10}+\frac9{10}-1=\frac35.
\]
On the intersection, the connected interval contains both targets and has length at least their separation. Thus
\[
\mathbb E_{P_{h_0}^{\mathrm{dir},-},n}\operatorname{length}(I)
\ge2\kappa h_0^\beta
\mathsf E_{\sigma,h_0}^{-}
(\mathsf F_{h_0,I}^{+}\cap\mathsf F_{h_0,I}^{-})
\ge\frac65\kappa h_0^\beta.
\]
Taking the infimum over honest intervals yields
\[
\inf_{I\in\mathcal I_{\beta,n,\sigma}}
\sup_{P\in\mathcal M_\beta(\sigma)}
\mathbb E_{P,n}\operatorname{length}(I)
\ge\frac65\kappa h_0^\beta.
\]
% lean: CausalSmith.Stat.RdTruesideNoiseFrontier.experiment_tv_le
% lean: CausalSmith.Stat.RdTruesideNoiseFrontier.twoPoint_absRisk_sum
% lean: CausalSmith.Stat.RdTruesideNoiseFrontier.twoPoint_interval_lower
% lean: CausalSmith.Stat.RdTruesideNoiseFrontier.direct_two_point_transfer

\item \textbf{Real lower bounds.}
For every \(P\in\mathcal M_\beta(\sigma)\), the mean bounds give
\[
|\theta(P)|
=|\mu_1(P,0)-\mu_0(P,0)|
\le\frac12.
\]
The constant-zero estimator consequently satisfies
\[
\sup_{P\in\mathcal M_\beta(\sigma)}
\mathbb E_{P,n}|0-\theta(P)|\le\frac12.
\]
The constant interval \([-1,1]\) satisfies
\[
\mathbb P_{P,n}\{\theta(P)\in[-1,1]\}=1,
\qquad
\mathbb E_{P,n}\operatorname{length}([-1,1])=2
\quad(P\in\mathcal M_\beta(\sigma)).
\]
It is therefore an admissible honest interval. These two procedures bound the extended minimax values by finite numbers, so conversion to real numbers preserves the preceding lower inequalities. Since \(h_0^\beta=r_{\mathrm{dir},n}\),
\[
R_\beta(n,\sigma)\ge\frac45\kappa r_{\mathrm{dir},n},
\qquad
\mathcal L_\beta(n,\sigma)\ge\frac65\kappa r_{\mathrm{dir},n}.
\]
Define
\[
c:=\frac45\kappa>0.
\]
As \(\kappa r_{\mathrm{dir},n}\ge0\), both criteria are at least
\(c\,r_{\mathrm{dir},n}\), uniformly over \(n\ge2\) and \(\sigma\in[0,1]\).
% lean: CausalSmith.Stat.RdTruesideNoiseFrontier.risk_value_ne_top
% lean: CausalSmith.Stat.RdTruesideNoiseFrontier.length_value_ne_top
% lean: CausalSmith.Stat.RdTruesideNoiseFrontier.direct_uniform_lower

\item \textbf{Attainment and minimax upper bounds.}
Set
\[
C:=2K_{\mathrm{cert}}>0,
\qquad
P_{\mathrm{unit}}:=P_1^{\mathrm{dir},+}.
\]
The witness verification above, at \(h=1\), gives
\[
P_{\mathrm{unit}}\in\mathcal M_\beta(0).
\]
For every \(n\ge2\), apply \cref{obj:thm:finite-certificate} with
\(\sigma=0\), \(L=1\), and \(P=P_{\mathrm{unit}}\). Its selected-procedure conclusions give
\[
I_{L_\star}\in\mathcal I_{\beta,n,0},
\]
\[
\sup_{P\in\mathcal M_\beta(0)}
\mathbb E_{P,n}
|\widehat\theta_{L_\star}-\theta(P)|
\le E_\beta(n,0),
\qquad
\sup_{P\in\mathcal M_\beta(0)}
\mathbb E_{P,n}\operatorname{length}(I_{L_\star})
\le2E_\beta(n,0).
\]
Since \(K_{\mathrm{cert}}r_{\mathrm{dir},n}\ge0\), the radius bound from the first step implies
\[
\sup_{P\in\mathcal M_\beta(0)}
\mathbb E_{P,n}
|\widehat\theta_{L_\star}-\theta(P)|
\le K_{\mathrm{cert}}r_{\mathrm{dir},n}
\le C r_{\mathrm{dir},n},
\]
\[
\sup_{P\in\mathcal M_\beta(0)}
\mathbb E_{P,n}\operatorname{length}(I_{L_\star})
\le2K_{\mathrm{cert}}r_{\mathrm{dir},n}
=C r_{\mathrm{dir},n}.
\]
Each extended minimax value is bounded above by the worst-case loss of its displayed admissible procedure. The selected interval's honesty supplies its admissibility for the length minimization. The finite upper bounds permit order-preserving conversion to real numbers, giving
\[
R_\beta(n,0)\le C r_{\mathrm{dir},n},
\qquad
\mathcal L_\beta(n,0)\le C r_{\mathrm{dir},n}.
\]
The constants \(c\) and \(C\) depend only on \(\beta\), and the displayed inequalities and honesty statement hold for every integer \(n\ge2\), as required.
% lean: CausalSmith.Stat.RdTruesideNoiseFrontier.risk_le_of_attainerRisk_le
% lean: CausalSmith.Stat.RdTruesideNoiseFrontier.lengthRisk_le_of_attainerLength_le
% lean: CausalSmith.Stat.RdTruesideNoiseFrontier.direct_reduction
\end{enumerate}
\end{proof}

\begin{proof}[Proof of \cref{obj:thm:uniform-frontier-resolution}]
\begin{enumerate}
\item Apply \cref{obj:thm:uniform-frontier} with the given \(\beta\). It supplies positive constants \(c,C\), depending only on \(\beta\), and a unique resolution in \([h_0,1]\). Its defining equation is equivalent to the equation in the present statement:
\[
\log n+\nu\log h_\star=\Psi(h_\star,\sigma)
\quad\Longleftrightarrow\quad
\nu\log(1/h_\star)+\Psi(h_\star,\sigma)=\log n.
\]
The rate and selected procedures are precisely those of
\cref{obj:def:uniform-rate,obj:def:public-selection}. Consequently, the two estimation and expected-length comparison chains, together with interval honesty, follow directly from \cref{obj:thm:uniform-frontier}. Its decision-class conclusions include procedures with an independent uniform seed. The same constants apply simultaneously for all \(n\ge2\) and \(\sigma\in[0,1]\), and hence along every shrinking-noise sequence.
% lean: CausalSmith.Stat.RdTruesideNoiseFrontier.uniform_frontier

\item The direct and intermediate formulas and the strict compact formula are supplied by \cref{obj:thm:uniform-frontier}. At the intermediate--compact equality, that result gives
\[
h_\star=\sigma^4,\qquad z=\sigma^{-2},\qquad \tau=\sigma^{-2},
\]
with both scalar equations satisfied. Thus the compact representation also includes the equality case. Its scalar equation uniquely determines \(\tau\) on the stated domain: for positive \(\sigma\), its left side has derivative
\[
\frac{2\nu}{\tau}+2+\log(\sigma^2\tau)>0
\qquad(\tau\ge\sigma^{-2}).
\]
The direct interface is likewise supplied by the same result:
\[
h_\star=\sigma=h_0.
\]
Finally, its specialization to \(\beta=1\) and \(\sigma=n^{-1/6}\) gives the intermediate branch for every \(n\ge2\) and the asserted comparison
\[
r_1(n,n^{-1/6})
\asymp
\frac{n^{-1/6}}{[1+\tfrac12\log n]^{3/2}}.
\]
These conclusions use the observed side label \(D=\mathbf1\{X\ge0\}\).
% lean: CausalSmith.Stat.RdTruesideNoiseFrontier.uniform_frontier_resolution
\end{enumerate}
\end{proof}
